\section{Appendix}

\begin{center}
\fbox{\begin{minipage}{0.86\textwidth}
\textbf{\color{journalgreen} Reproducibility}\par
Run from the project root after configuring local data paths and credentials where
needed:\par
\vspace{2pt}
\texttt{make option-paper}\par
\texttt{make paper}\par
\texttt{make verify}\par
\texttt{make test}\par
\end{minipage}}
\end{center}

\subsection{Exposures, attribution, and realized-path diagnostics}

This first group records what the strategies actually held and where the P\&L came from.
Figure~\ref{fig:growth_all_app} keeps every generated strategy visible on one growth
panel, Figure~\ref{fig:risk_calibration_app} compares predicted and realized volatility,
Tables~\ref{tab:premium_exposure_app} and~\ref{tab:greek_exposure_app} record premium and
Greek exposures, and Table~\ref{tab:pnl_attribution_app} decomposes mean P\&L into Greek
and residual components.

\begin{figure}[H]
\centering
\IfFileExists{figures/portfolio_growth_all_strategies.pdf}
  {\includegraphics[width=0.86\textwidth]{portfolio_growth_all_strategies.pdf}}
  {\fbox{\parbox{0.78\textwidth}{Run the empirical pipeline to generate the all-strategy growth diagnostic.}}}
\caption{All-strategy out-of-sample growth diagnostic. The main text shows the central
strategies; this appendix figure keeps every generated strategy visible for auditability.}
\label{fig:growth_all_app}
\end{figure}

\begin{figure}[H]
\centering
\IfFileExists{figures/risk_calibration.pdf}
  {\includegraphics[width=0.72\textwidth]{risk_calibration.pdf}}
  {\fbox{\parbox{0.75\textwidth}{Run the empirical pipeline to generate the risk calibration figure.}}}
\caption{Predicted versus realized monthly volatility by option strategy.}
\label{fig:risk_calibration_app}
\end{figure}

\begin{table}[H]
\centering
\scriptsize
\IfFileExists{tables/exposure_summary.tex}
  {\resizebox{0.92\textwidth}{!}{\begin{tabular}{llrrrrrr}
\toprule
Strategy & Book & Gross NAV & Net NAV & Long premium paid & Short premium sold & Net option premium & Short gross \\
\midrule
Equity-option Greek Markowitz & Options & 1.000 & 0.500 & 0.750 & 0.250 & 0.500 & 0.250 \\
Greek Markowitz + VIX & Options & 1.000 & 0.500 & 0.750 & 0.250 & 0.500 & 0.250 \\
Beta/delta-neutral + VIX & Options & 1.000 & 0.500 & 0.750 & 0.250 & 0.500 & 0.250 \\
Equal premium & Options & 1.000 & 1.000 & 1.000 & 0.000 & 1.000 & 0.000 \\
Equal risk & Options & 1.000 & 1.000 & 1.000 & 0.000 & 1.000 & 0.000 \\
VIX hedge sleeve & Options & 1.000 & 1.000 & 1.000 & 0.000 & 1.000 & 0.000 \\
Delta-matched equities & Equities & 3.307 & 3.307 &  &  &  & 0.000 \\
Underlying Markowitz & Equities & 1.000 & 0.539 &  &  &  & 0.231 \\
\bottomrule
\end{tabular}
}}
  {\begin{tabular}{llll}\toprule Strategy & Gross NAV & Net NAV & Short gross \\\midrule Missing & -- & -- & -- \\\bottomrule\end{tabular}}
\caption{Premium and sign exposures. The table records how much option premium is paid,
sold, and netted before the strategy return is realized.}
\label{tab:premium_exposure_app}
\end{table}

\begin{table}[H]
\centering
\scriptsize
\IfFileExists{tables/greek_exposure_summary.tex}
  {\resizebox{\textwidth}{!}{\begin{tabular}{llrrrrrrrrrrr}
\toprule
Strategy & Book & Net delta & Gross delta & Net gamma & Net vega & Equity vega & VIX vega & Beta SPY proxy & Worst stress return & VIX option gross & Call gross & Put gross \\
\midrule
Equity-option Greek Markowitz & Options & 2.336 & 4.055 & 5.376 & 0.004 & 0.004 & 0.000 & 3.000 & -0.207 & 0.000 & 0.682 & 0.318 \\
Greek Markowitz + VIX & Options & 3.307 & 3.307 & 5.029 & -0.108 & 0.005 & -0.113 & 2.342 & -0.350 & 0.200 & 0.633 & 0.367 \\
Beta/delta-neutral + VIX & Options & -0.086 & 3.893 & 2.186 & -0.099 & 0.005 & -0.104 & -0.145 & -0.010 & 0.200 & 0.588 & 0.412 \\
Cost-aware Sortino + VIX & Options & 6.288 & 16.065 & 129.870 & -0.652 & 0.035 & -0.687 & 3.000 & -0.276 & 0.200 & 0.679 & 0.321 \\
Equal premium & Options & 1.208 & 2.202 & 359.583 & 0.211 & 0.041 & 0.170 & 0.897 & 0.258 & 0.093 & 0.537 & 0.463 \\
Equal risk & Options & 0.650 & 1.988 & 175.662 & 0.116 & 0.027 & 0.089 & 0.689 & 0.126 & 0.065 & 0.561 & 0.439 \\
VIX hedge sleeve & Options & 0.000 & 0.000 & 18.435 & 1.572 & 0.000 & 1.572 & -23.137 & -0.406 & 1.000 & 1.000 & 0.000 \\
Delta-matched equities & Equities & 3.307 & 3.307 & 0.000 & 0.000 & 0.000 & 0.000 &  &  & 0.000 &  &  \\
Underlying Markowitz & Equities & 0.539 & 1.000 & 0.000 & 0.000 & 0.000 & 0.000 &  &  & 0.000 &  &  \\
\bottomrule
\end{tabular}
}}
  {\begin{tabular}{llll}\toprule Strategy & Delta & Vega & Stress \\\midrule Missing & -- & -- & -- \\\bottomrule\end{tabular}}
\caption{Greek, beta-proxy, VIX-vega, stress, and option-right exposures.}
\label{tab:greek_exposure_app}
\end{table}

\begin{table}[H]
\centering
\scriptsize
\IfFileExists{tables/pnl_attribution.tex}
  {\resizebox{\textwidth}{!}{\begin{tabular}{lrrrrrrrrrrrr}
\toprule
Strategy & Equity delta & Equity gamma & Equity vega & VIX-forward delta & VIX-forward gamma & VIX-option vega & Theta/carry & VX roll & Skew/tail & Residual & Realized ann. mean & Residual vol share \\
\midrule
Equity-option Greek Markowitz & 0.521 & 0.898 & 0.002 & 0.000 & 0.000 & 0.000 & -0.003 & 0.000 & 0.000 & -0.524 & 0.894 & 0.570 \\
Greek Markowitz + VIX & 0.489 & 1.475 & 0.001 & 0.767 & -0.867 & -0.227 & 1.057 & 0.000 & 0.000 & -1.447 & 1.246 & 0.950 \\
Beta/delta-neutral + VIX & 0.168 & 1.486 & 0.001 & 0.554 & -0.745 & -0.190 & 1.034 & 0.000 & 0.000 & -1.135 & 1.172 & 1.069 \\
Cost-aware Sortino + VIX & 2.064 & 8.577 & 0.004 & 0.285 & -1.134 & -0.112 & 2.169 & 0.000 & 0.000 & -6.833 & 5.020 & 0.710 \\
Equal premium & 0.062 & 7.928 & 0.006 & -0.016 & 0.358 & 0.045 & -1.016 & 0.000 & 0.000 & -7.423 & -0.056 & 0.354 \\
Equal risk & 0.027 & 5.625 & 0.004 & -0.096 & 0.270 & 0.049 & -0.592 & 0.000 & 0.000 & -5.184 & 0.104 & 0.514 \\
VIX hedge sleeve & 0.000 & 0.000 & 0.000 & 2.083 & 2.615 & 0.135 & -12.254 & 0.000 & 0.000 & -1.865 & -9.285 & 4.858 \\
\bottomrule
\end{tabular}
}}
  {\begin{tabular}{llll}\toprule Strategy & Delta & Vega & Residual \\\midrule Missing & -- & -- & -- \\\bottomrule\end{tabular}}
\caption{Annualized mean P\&L attribution. Equity-option Greeks and VIX-forward Greeks
are separated; components reconcile to realized mean after the residual term.}
\label{tab:pnl_attribution_app}
\end{table}

\subsection{Costs, capacity, and settlement diagnostics}

The second group makes timing, settlement, and cost conventions visible before any
post-cost number is interpreted. Table~\ref{tab:timing_app} records return-construction
timing, Table~\ref{tab:trading_audit_app} audits the point-in-time holding model, and
Table~\ref{tab:vix_download_audit_app} documents which required VIX expiries have parsed
exact VRO/SOQ settlement. The expiry-level exact-versus-proxy settlement comparison spans
110 expiries and is too long to typeset here; it ships as the machine-readable artifact
\path{artifacts/vix_settlement_audit.csv} (mirrored in
\path{tables/vix_settlement_audit.tex}). Table~\ref{tab:costs_app} itemizes the
conservative stress-cost layer, Table~\ref{tab:capacity_market_impact_app} reports the
executable-cost scenarios, and Tables~\ref{tab:performance_diag_app}
and~\ref{tab:performance_net_diag_app} give the gross and stress-cost NAV diagnostics,
including the bootstrap Sharpe confidence intervals quoted in Section~\ref{sec:results}.

\begin{table}[H]
\centering
\scriptsize
\IfFileExists{tables/timing_diagnostics.tex}
  {\resizebox{0.92\textwidth}{!}{\input{tables/timing_diagnostics}}}
  {\begin{tabular}{ll}\toprule Diagnostic & Value \\\midrule Missing & Run the empirical pipeline \\\bottomrule\end{tabular}}
\caption{Timing and return-construction diagnostics.}
\label{tab:timing_app}
\end{table}

\begin{table}[H]
\centering
\scriptsize
\IfFileExists{tables/trading_data_audit.tex}
  {\resizebox{\textwidth}{!}{\begin{tabular}{lll}
\toprule
Check & Value & Pass \\
\midrule
OPRA-derived rows after liquidity filters & 160,315 & yes \\
Representative option choices & 5,825 & yes \\
Nonmissing option return cells & 5,777 & yes \\
Max OPRA spot vs raw close mismatch & 2.129e-16 & yes \\
All option choices made before payoff date & True & yes \\
All payoff dates no later than next rebalance & True & yes \\
Exact listed-expiry close share & 1.000 & yes \\
Max expiry-to-payoff-date lag in days & 0 & yes \\
Holding period days min/median/max & 15/18.0/21 & yes \\
Rows requiring split conversion during holding & 13 & yes \\
Detected raw-close split/unit events & 7 & yes \\
Minimum single-option return & -1.000 & yes \\
Maximum single-option return & 48.700 & yes \\
Max training decision before first test decision & 2020-11-30 before 2020-12-31 & yes \\
Raw VIX option rows after contract parsing & 103,650 & yes \\
Representative VIX option choices & 630 & yes \\
VIX option holding ledger rows & 536 & yes \\
Duplicate VIX date-symbol rows after dedupe & 0 & yes \\
Black-76 VX-forward Greek coverage & 1.000 & yes \\
VIX settlement source & vro\_soq\_exact:536 & yes \\
All VIX decisions before payoff dates & True & yes \\
Minimum VIX long-option return & -1.000 & yes \\
\bottomrule
\end{tabular}
}}
  {\begin{tabular}{lll}\toprule Check & Value & Pass \\\midrule Missing & -- & -- \\\bottomrule\end{tabular}}
\caption{Trading-data audit for the point-in-time holding model. Rows include VIX option
settlement-source flags; proxy rows, if present, remain diagnostic.}
\label{tab:trading_audit_app}
\end{table}

\begin{table}[H]
\centering
\scriptsize
\IfFileExists{tables/vix_required_settlement_download_audit.tex}
  {\resizebox{0.86\textwidth}{!}{\input{tables/vix_required_settlement_download_audit}}}
  {\begin{tabular}{llll}\toprule Year & Exact & Component & Expiries \\\midrule Missing & -- & -- & -- \\\bottomrule\end{tabular}}
\caption{Public Cboe VRO/SOQ download audit for required VIX expiries. Rows with
\texttt{Exact scalar parsed = False} are not used in settlement-gated VIX option returns;
the excluded expiries cluster in 2015--2018 (Section~\ref{sec:data} discusses the resulting coverage
bias).}
\label{tab:vix_download_audit_app}
\end{table}

\begin{table}[H]
\centering
\scriptsize
\IfFileExists{tables/cost_capacity_margin_diagnostics.tex}
  {\resizebox{\textwidth}{!}{\begin{tabular}{lrrrrrrrr}
\toprule
Strategy & Mean monthly cost & Ann. cost drag & Spread share & Fee share & Mean margin/NAV & Mean stress margin/NAV & Mean assignment notional/NAV & Capacity penalized share \\
\midrule
Beta/delta-neutral + VIX & 0.422 & 5.068 & 0.099 & 0.012 & 4.029 & 0.235 & 0.000 & 0.273 \\
Equal premium & 0.259 & 3.109 & 0.115 & 0.025 & 0.953 & 0.238 & 0.000 & 0.381 \\
Equal risk & 0.115 & 1.381 & 0.213 & 0.035 & 0.920 & 0.230 & 0.000 & 0.347 \\
Equity-option Greek Markowitz & 0.733 & 8.794 & 0.021 & 0.003 & 4.610 & 0.224 & 0.039 & 0.335 \\
Greek Markowitz + VIX & 0.610 & 7.325 & 0.069 & 0.008 & 3.526 & 0.234 & 0.000 & 0.290 \\
VIX hedge sleeve & 0.430 & 5.165 & 0.347 & 0.046 & 0.978 & 0.244 & 0.000 & 0.847 \\
\bottomrule
\end{tabular}
}}
  {\begin{tabular}{llll}\toprule Strategy & Cost & Capacity & Margin \\\midrule Missing & -- & -- & -- \\\bottomrule\end{tabular}}
\caption{Conservative stress-cost layer diagnostics. Costs include round-trip spread
crossing, explicit fees, slippage, borrow, capacity penalties, simulated margin funding,
and assignment-risk penalties where flagged; this is the layer behind the ``Post-cost
research'' rows of Table~\ref{tab:performance}, not the executable-cost scenarios of
Table~\ref{tab:capacity_market_impact_app}. These are not live broker fills.}
\label{tab:costs_app}
\end{table}

\begin{table}[H]
\centering
\scriptsize
\IfFileExists{tables/capacity_market_impact_diagnostics.tex}
  {\resizebox{\textwidth}{!}{\begin{tabular}{lllllllll}
\toprule
Strategy & Scenario & Avg contracts traded & Avg quoted spread & Avg monthly cost & Max capacity used & Rejected trades & Mean margin/NAV & Mean stress capital/NAV \\
\midrule
Beta/delta-neutral + VIX & full\_spread & 56.74896912971811 & 0.03509309915904966 & 0.03198047526748731 & 0.9997072467259026 & 867 & 1.6454288980146308 & 0.11046921803293545 \\
Beta/delta-neutral + VIX & half\_spread & 56.74896912971811 & 0.03509309915904966 & 0.018746121851315658 & 0.9997072467259026 & 867 & 1.6454288980146308 & 0.11046921803293545 \\
Beta/delta-neutral + VIX & mid & 56.74896912971811 & 0.03509309915904966 & 0.005511768435143996 & 0.9997072467259026 & 867 & 1.6454288980146308 & 0.11046921803293545 \\
Cost-aware Sortino + VIX & full\_spread & 486.94649539475853 & 0.0764124240818374 & 0.015719557112203447 & 0.9984987277951501 & 265 & 0.6759167177588942 & 0.07023236901345059 \\
Cost-aware Sortino + VIX & half\_spread & 486.94649539475853 & 0.0764124240818374 & 0.00972463183676715 & 0.9984987277951501 & 265 & 0.6759167177588942 & 0.07023236901345059 \\
Cost-aware Sortino + VIX & mid & 486.94649539475853 & 0.0764124240818374 & 0.003729706561330857 & 0.9984987277951501 & 265 & 0.6759167177588942 & 0.07023236901345059 \\
Equal premium & full\_spread & 65.36655726350517 & 0.03450510899698828 & 0.027393513458707507 & 0.9993803841618195 & 1330 & 0.5895061728395058 & 0.14737654320987645 \\
Equal premium & half\_spread & 65.36655726350517 & 0.03450510899698828 & 0.017223026084595215 & 0.9993803841618195 & 1330 & 0.5895061728395058 & 0.14737654320987645 \\
Equal premium & mid & 65.36655726350517 & 0.03450510899698828 & 0.007052538710482932 & 0.9993803841618195 & 1330 & 0.5895061728395058 & 0.14737654320987645 \\
Equal risk & full\_spread & 43.657154853888706 & 0.03490292057829765 & 0.021254428248565684 & 0.9990134907217686 & 1222 & 0.5170241047560368 & 0.1292560261890092 \\
Equal risk & half\_spread & 43.657154853888706 & 0.03490292057829765 & 0.013262882446192524 & 0.9990134907217686 & 1222 & 0.5170241047560368 & 0.1292560261890092 \\
Equal risk & mid & 43.657154853888706 & 0.03490292057829765 & 0.005271336643819363 & 0.9990134907217686 & 1222 & 0.5170241047560368 & 0.1292560261890092 \\
Equity-option Greek Markowitz & full\_spread & 18.797109490588205 & 0.017435798276330978 & 0.006679109598942776 & 0.9901002345178692 & 957 & 1.3950575743846914 & 0.07960287865721945 \\
Equity-option Greek Markowitz & half\_spread & 18.797109490588205 & 0.017435798276330978 & 0.0044518190699674685 & 0.9901002345178692 & 957 & 1.3950575743846914 & 0.07960287865721945 \\
Equity-option Greek Markowitz & mid & 18.797109490588205 & 0.017435798276330978 & 0.002224528540992161 & 0.9901002345178692 & 957 & 1.3950575743846914 & 0.07960287865721945 \\
Greek Markowitz + VIX & full\_spread & 49.294674731990554 & 0.035957144195152634 & 0.026579996078377945 & 0.99884496167564 & 898 & 1.5123553671138326 & 0.1014995147957134 \\
Greek Markowitz + VIX & half\_spread & 49.294674731990554 & 0.035957144195152634 & 0.01563456322750504 & 0.99884496167564 & 898 & 1.5123553671138326 & 0.1014995147957134 \\
Greek Markowitz + VIX & mid & 49.294674731990554 & 0.035957144195152634 & 0.0046891303766321375 & 0.99884496167564 & 898 & 1.5123553671138326 & 0.1014995147957134 \\
VIX hedge sleeve & full\_spread & 2410.1807707861826 & 0.15 & 0.07580327078625228 & 0.9885078059495911 & 153 & 0.14999999999999997 & 0.03749999999999999 \\
VIX hedge sleeve & half\_spread & 2410.1807707861826 & 0.15 & 0.042053270786252275 & 0.9885078059495911 & 153 & 0.14999999999999997 & 0.03749999999999999 \\
VIX hedge sleeve & mid & 2410.1807707861826 & 0.15 & 0.008303270786252284 & 0.9885078059495911 & 153 & 0.14999999999999997 & 0.03749999999999999 \\
\bottomrule
\end{tabular}
}}
  {\begin{tabular}{llll}\toprule Strategy & Scenario & Contracts & Capacity \\\midrule Missing & -- & -- & -- \\\bottomrule\end{tabular}}
\caption{Executable-cost, capacity, and market-impact diagnostics by scenario
(entry-friction layer). Rows are generated from the cost-scenario ledger and record
no-fill/capacity failures rather than assuming midpoint execution.}
\label{tab:capacity_market_impact_app}
\end{table}

The CBBO surface is an optional research-cost refinement: panel CBBO spreads take
precedence, then exact decision-date surface spreads, then class-default spreads. The
coverage table records which source supplied the relative spread by asset class.

\begin{table}[H]
\centering
\scriptsize
\IfFileExists{tables/cost_input_spread_source_coverage.tex}
  {\resizebox{0.84\textwidth}{!}{\begin{tabular}{llrr}
\toprule
Spread source & Asset class & Rows & Mean relative spread \\
\midrule
default & VIX & 536 & 0.150 \\
panel\_cbbo & equity & 5777 & 0.026 \\
\bottomrule
\end{tabular}
}}
  {\begin{tabular}{llll}\toprule Source & Asset class & Rows & Spread \\\midrule Missing & -- & -- & -- \\\bottomrule\end{tabular}}
\caption{Cost-input spread-source coverage. Rows distinguish direct panel CBBO spreads,
derived CBBO surface spreads, and class-default spread assumptions.}
\label{tab:cost_input_spread_source_app}
\end{table}

The repaired execution scenarios are execution-sensitivity diagnostics. They use the same
decision-date cost-input row that fired the original gate; missing cost input and
assignment or dividend hard gates are not repaired.

\begin{table}[H]
\centering
\scriptsize
\IfFileExists{tables/execution_repair_diagnostics.tex}
  {\resizebox{0.86\textwidth}{!}{\begin{tabular}{lrrrr}
\toprule
Scenario & Repaired orders & Partial fills & Avg extra cost per repaired order & Foregone gross from unfilled remainders \\
\midrule
full\_spread\_repaired & 4181 & 4181 & 0.000 & -4.580 \\
half\_spread\_repaired & 4181 & 4181 & 0.000 & -4.580 \\
mid\_repaired & 4181 & 4181 & 0.000 & -4.580 \\
\bottomrule
\end{tabular}
}}
  {\begin{tabular}{llll}\toprule Scenario & Repairs & Partial fills & Cost \\\midrule Missing & -- & -- & -- \\\bottomrule\end{tabular}}
\caption{Execution-repair diagnostics for \texttt{\_repaired} scenarios. Counts and
costs are diagnostics only and are not headline performance results.}
\label{tab:execution_repair_diagnostics_app}
\end{table}

The comparison table reports how repaired execution-sensitivity scenarios differ from the
fail-closed cost scenarios. It is a sensitivity check around the execution gates, not a
claim that the repaired fills are broker-executed.

\begin{table}[H]
\centering
\scriptsize
\IfFileExists{tables/execution_repair_comparison.tex}
  {\resizebox{\textwidth}{!}{\begin{tabular}{llrrrrr}
\toprule
Strategy & Scenario & Fail-closed Sharpe & Repaired Sharpe & Fail-closed ann. return & Repaired ann. return & Repair uplift (ann.) \\
\midrule
Equity-option Greek Markowitz & mid & 1.024 & 1.195 & 0.488 & 0.694 & 0.206 \\
Greek Markowitz + VIX & mid & 0.771 & 1.355 & 0.424 & 0.823 & 0.400 \\
Beta/delta-neutral + VIX & mid & 1.265 & 1.622 & 0.691 & 0.915 & 0.224 \\
Cost-aware Sortino + VIX & mid & 0.142 & 0.592 & 0.214 & 1.424 & 1.211 \\
Equal premium & mid & 0.052 & 0.017 & 0.061 & 0.025 & -0.037 \\
Equal risk & mid & 0.137 & 0.143 & 0.108 & 0.139 & 0.031 \\
VIX hedge sleeve & mid & -1.700 & -4.256 & -1.417 & -4.277 & -2.860 \\
Equity-option Greek Markowitz & half\_spread & 0.969 & 1.122 & 0.462 & 0.652 & 0.191 \\
Greek Markowitz + VIX & half\_spread & 0.534 & 1.051 & 0.293 & 0.639 & 0.347 \\
Beta/delta-neutral + VIX & half\_spread & 0.971 & 1.275 & 0.532 & 0.721 & 0.189 \\
Cost-aware Sortino + VIX & half\_spread & 0.094 & 0.540 & 0.142 & 1.298 & 1.156 \\
Equal premium & half\_spread & -0.052 & -0.085 & -0.061 & -0.119 & -0.058 \\
Equal risk & half\_spread & 0.015 & 0.024 & 0.012 & 0.024 & 0.012 \\
VIX hedge sleeve & half\_spread & -1.742 & -4.458 & -1.552 & -4.675 & -3.123 \\
Equity-option Greek Markowitz & full\_spread & 0.913 & 1.049 & 0.435 & 0.610 & 0.175 \\
Greek Markowitz + VIX & full\_spread & 0.294 & 0.748 & 0.161 & 0.455 & 0.294 \\
Beta/delta-neutral + VIX & full\_spread & 0.678 & 0.929 & 0.373 & 0.527 & 0.153 \\
Cost-aware Sortino + VIX & full\_spread & 0.047 & 0.487 & 0.070 & 1.172 & 1.102 \\
Equal premium & full\_spread & -0.156 & -0.188 & -0.183 & -0.263 & -0.080 \\
Equal risk & full\_spread & -0.107 & -0.094 & -0.084 & -0.091 & -0.008 \\
VIX hedge sleeve & full\_spread & -1.777 & -4.635 & -1.687 & -5.072 & -3.385 \\
\bottomrule
\end{tabular}
}}
  {\begin{tabular}{llll}\toprule Strategy & Scenario & Fail closed & Repaired \\\midrule Missing & -- & -- & -- \\\bottomrule\end{tabular}}
\caption{Fail-closed versus repaired execution-sensitivity performance. Repaired columns
remain outside the headline growth exhibits and the reality-check family.}
\label{tab:execution_repair_comparison_app}
\end{table}

\clearpage

The cost-aware Sortino variant uses train-window entry-cost estimates before solving the
Sortino objective. It is included as a research diagnostic and is not part of the headline
strategy set or tail-path simulation set.

\begin{table}[H]
\centering
\scriptsize
\IfFileExists{tables/sortino_objective_diagnostics.tex}
  {\resizebox{0.78\textwidth}{!}{\begin{tabular}{lllrrrrrrrrr}
\toprule
Status & Method & Objective & Net mean & Gross mean & Entry cost & Downside deviation & Net Sortino & Target & Scenarios & Downside-free & Train scenarios \\
\midrule
optimal & linprog\_net\_mean\_downside\_free & max\_sortino\_net\_of\_costs & 0.147 & 0.184 & 0.037 & 0.357 & 0.412 & 0.000 & 69 & True & 69 \\
\bottomrule
\end{tabular}
}}
  {\begin{tabular}{llll}\toprule Objective & Status & Value & Assets \\\midrule Missing & -- & -- & -- \\\bottomrule\end{tabular}}
\caption{Cost-aware Sortino objective diagnostics. Entry costs are estimated from the
training window and imputed from class defaults when an asset has no train-window cost
observation.}
\label{tab:sortino_objective_diagnostics_app}
\end{table}

\clearpage

\begin{table}[H]
\centering
\scriptsize
\IfFileExists{tables/portfolio_performance_diagnostics.tex}
  {\resizebox{\textwidth}{!}{\begin{tabular}{lrrrrrrrr}
\toprule
Strategy & Downside ann. dev & Max drawdown & Worst month & SR 90\% CI lo & SR 90\% CI hi & Pred./realized vol & Gross NAV & Net NAV \\
\midrule
Equity-option Greek Markowitz & 0.505 & -0.876 & -0.455 & 0.842 & 0.083 & 1.089 & 1.000 & 0.500 \\
Greek Markowitz + VIX & 0.368 & -0.613 & -0.371 & 1.374 & 0.602 & 1.385 & 1.000 & 0.500 \\
Beta/delta-neutral + VIX & 0.344 & -0.570 & -0.354 & 1.414 & 0.676 & 1.478 & 1.000 & 0.500 \\
Equal premium & 1.004 & -1.000 & -0.662 & -0.030 & -0.635 & 2.782 & 1.000 & 1.000 \\
Equal risk & 0.713 & -0.976 & -0.516 & 0.082 & -0.558 & 1.814 & 1.000 & 1.000 \\
VIX hedge sleeve & 3.006 &  & -1.000 & -6.429 & -9.004 & 11.150 & 1.000 & 1.000 \\
Delta-matched equities & 0.473 & -0.894 & -0.454 & 1.176 & 0.372 &  & 3.307 & 3.307 \\
Underlying Markowitz & 0.096 & -0.303 & -0.073 & 0.909 & 0.103 &  & 1.000 & 0.539 \\
\bottomrule
\end{tabular}
}}
  {\begin{tabular}{llll}\toprule Strategy & Downside & Drawdown & NAV \\\midrule Missing & -- & -- & -- \\\bottomrule\end{tabular}}
\caption{Gross performance diagnostics and NAV accounting. The VIX hedge sleeve's maximum
drawdown is undefined (``--'') because its compounded wealth path is absorbed at zero
(Section~\ref{sec:metrics}).}
\label{tab:performance_diag_app}
\end{table}

\begin{table}[H]
\centering
\scriptsize
\IfFileExists{tables/portfolio_performance_net_diagnostics.tex}
  {\resizebox{\textwidth}{!}{\begin{tabular}{lrrrrrrrr}
\toprule
Strategy & Downside ann. dev & Max drawdown & Worst month & SR 90\% CI lo & SR 90\% CI hi & Pred./realized vol & Gross NAV & Net NAV \\
\midrule
Equity-option Greek Markowitz & 5.864 & -1.243 & -10.800 & -2.209 & -1.164 &  & 1.000 & 0.500 \\
Greek Markowitz + VIX & 5.074 & -1.183 & -9.226 & -1.871 & -0.994 &  & 1.000 & 0.500 \\
Beta/delta-neutral + VIX & 2.997 & -1.000 & -5.540 & -1.961 & -1.094 &  & 1.000 & 0.500 \\
Cost-aware Sortino + VIX & 154.544 & -79.319 & -320.564 & -2.019 & -0.961 &  & 1.000 & 0.600 \\
Equal premium & 2.107 & -1.037 & -2.431 & -1.975 & -0.816 &  & 1.000 & 1.000 \\
Equal risk & 1.065 & -1.000 & -0.974 & -1.576 & -0.486 &  & 1.000 & 1.000 \\
VIX hedge sleeve & 5.876 & -1.159 & -9.482 & -7.714 & -2.676 &  & 1.000 & 1.000 \\
Delta-matched equities & 0.473 & -0.894 & -0.454 & 0.372 & 2.182 &  &  &  \\
Underlying Markowitz & 0.096 & -0.303 & -0.073 & 0.103 & 1.812 &  &  &  \\
\bottomrule
\end{tabular}
}}
  {\begin{tabular}{llll}\toprule Strategy & Downside & Drawdown & NAV \\\midrule Missing & -- & -- & -- \\\bottomrule\end{tabular}}
\caption{Stress-cost performance diagnostics and NAV accounting.}
\label{tab:performance_net_diag_app}
\end{table}

\subsection{Inference and multiple testing}

The third group carries the statistical evidence discussed in Section~\ref{sec:results}.
Table~\ref{tab:inference_app} is the full bootstrap interval ledger,
Table~\ref{tab:reality_check_app} reports per-variant probabilistic and deflated Sharpe
ratios with the centered reality-check family statistics, and
Table~\ref{tab:factor_regression_full_app} is the complete factor-regression ledger behind
the headline regression table.

\begin{table}[H]
\centering
\scriptsize
\IfFileExists{tables/inference_summary.tex}
  {\resizebox{0.92\textwidth}{!}{\begin{tabular}{llrrrrlrrl}
\toprule
Strategy & Metric & Estimate & CI lo & CI hi & N & Method & Block size & Seed & Return basis \\
\midrule
Equity-option Greek Markowitz & Sharpe & 0.842 & 0.083 & 1.679 & 60 & circular block bootstrap & 8 & 20260625 & Gross before costs \\
Equity-option Greek Markowitz & Sortino & 1.770 & 0.141 & 4.581 & 60 & circular block bootstrap & 8 & 20260625 & Gross before costs \\
Equity-option Greek Markowitz & Calmar & 1.020 & 0.083 & 2.828 & 60 & circular block bootstrap & 8 & 20260625 & Gross before costs \\
Equity-option Greek Markowitz & Omega & 1.888 & 1.064 & 3.546 & 60 & circular block bootstrap & 8 & 20260625 & Gross before costs \\
Equity-option Greek Markowitz & Information Ratio & 0.598 & -0.430 & 1.571 & 60 & circular block bootstrap & 8 & 20260625 & Gross before costs \\
Greek Markowitz + VIX & Sharpe & 1.374 & 0.602 & 2.222 & 60 & circular block bootstrap & 8 & 20260625 & Gross before costs \\
Greek Markowitz + VIX & Sortino & 3.385 & 1.188 & 7.275 & 60 & circular block bootstrap & 8 & 20260625 & Gross before costs \\
Greek Markowitz + VIX & Calmar & 2.034 & 0.700 & 4.901 & 60 & circular block bootstrap & 8 & 20260625 & Gross before costs \\
Greek Markowitz + VIX & Omega & 2.784 & 1.562 & 5.251 & 60 & circular block bootstrap & 8 & 20260625 & Gross before costs \\
Greek Markowitz + VIX & Information Ratio & 1.200 & 0.347 & 2.129 & 60 & circular block bootstrap & 8 & 20260625 & Gross before costs \\
Beta/delta-neutral + VIX & Sharpe & 1.414 & 0.677 & 2.220 & 60 & circular block bootstrap & 8 & 20260625 & Gross before costs \\
Beta/delta-neutral + VIX & Sortino & 3.410 & 1.335 & 6.541 & 60 & circular block bootstrap & 8 & 20260625 & Gross before costs \\
Beta/delta-neutral + VIX & Calmar & 2.057 & 0.701 & 5.611 & 60 & circular block bootstrap & 8 & 20260625 & Gross before costs \\
Beta/delta-neutral + VIX & Omega & 2.948 & 1.693 & 5.333 & 60 & circular block bootstrap & 8 & 20260625 & Gross before costs \\
Beta/delta-neutral + VIX & Information Ratio & 1.147 & 0.332 & 1.962 & 60 & circular block bootstrap & 8 & 20260625 & Gross before costs \\
Cost-aware Sortino + VIX & Sharpe & 1.014 & 0.466 & 1.568 & 60 & circular block bootstrap & 8 & 20260625 & Gross before costs \\
Cost-aware Sortino + VIX & Sortino & 3.349 & 1.250 & 6.124 & 60 & circular block bootstrap & 8 & 20260625 & Gross before costs \\
Cost-aware Sortino + VIX & Calmar & 4.993 & 0.053 & 5.602 & 60 & circular block bootstrap & 8 & 20260625 & Gross before costs \\
Cost-aware Sortino + VIX & Omega & 2.471 & 1.510 & 4.021 & 60 & circular block bootstrap & 8 & 20260625 & Gross before costs \\
Cost-aware Sortino + VIX & Information Ratio & 0.971 & 0.373 & 1.518 & 60 & circular block bootstrap & 8 & 20260625 & Gross before costs \\
Equal premium & Sharpe & -0.030 & -0.635 & 0.437 & 60 & circular block bootstrap & 8 & 20260625 & Gross before costs \\
Equal premium & Sortino & -0.055 & -0.946 & 0.920 & 60 & circular block bootstrap & 8 & 20260625 & Gross before costs \\
Equal premium & Calmar & -0.056 & -1.051 & 0.849 & 60 & circular block bootstrap & 8 & 20260625 & Gross before costs \\
Equal premium & Omega & 0.977 & 0.624 & 1.416 & 60 & circular block bootstrap & 8 & 20260625 & Gross before costs \\
Equal premium & Information Ratio & -0.120 & -0.779 & 0.351 & 60 & circular block bootstrap & 8 & 20260625 & Gross before costs \\
Equal risk & Sharpe & 0.082 & -0.558 & 0.627 & 60 & circular block bootstrap & 8 & 20260625 & Gross before costs \\
Equal risk & Sortino & 0.146 & -0.821 & 1.305 & 60 & circular block bootstrap & 8 & 20260625 & Gross before costs \\
Equal risk & Calmar & 0.107 & -0.667 & 0.954 & 60 & circular block bootstrap & 8 & 20260625 & Gross before costs \\
Equal risk & Omega & 1.064 & 0.662 & 1.621 & 60 & circular block bootstrap & 8 & 20260625 & Gross before costs \\
Equal risk & Information Ratio & -0.058 & -0.732 & 0.485 & 60 & circular block bootstrap & 8 & 20260625 & Gross before costs \\
VIX hedge sleeve & Sharpe & -6.429 & -9.004 & -4.921 & 60 & circular block bootstrap & 8 & 20260625 & Gross before costs \\
VIX hedge sleeve & Sortino & -3.089 & -3.253 & -2.904 & 60 & circular block bootstrap & 8 & 20260625 & Gross before costs \\
VIX hedge sleeve & Calmar &  & -10.058 & -8.262 & 60 & circular block bootstrap & 8 & 20260625 & Gross before costs \\
VIX hedge sleeve & Omega & 0.040 & 0.016 & 0.071 & 60 & circular block bootstrap & 8 & 20260625 & Gross before costs \\
VIX hedge sleeve & Information Ratio & -6.132 & -8.858 & -4.691 & 60 & circular block bootstrap & 8 & 20260625 & Gross before costs \\
Delta-matched equities & Sharpe & 1.176 & 0.372 & 2.182 & 60 & circular block bootstrap & 8 & 20260625 & Gross before costs \\
Delta-matched equities & Sortino & 2.317 & 0.615 & 5.886 & 60 & circular block bootstrap & 8 & 20260625 & Gross before costs \\
Delta-matched equities & Calmar & 1.225 & 0.388 & 3.716 & 60 & circular block bootstrap & 8 & 20260625 & Gross before costs \\
Delta-matched equities & Omega & 2.367 & 1.301 & 5.182 & 60 & circular block bootstrap & 8 & 20260625 & Gross before costs \\
Delta-matched equities & Information Ratio & 1.041 & 0.189 & 2.072 & 60 & circular block bootstrap & 8 & 20260625 & Gross before costs \\
Underlying Markowitz & Sharpe & 0.909 & 0.103 & 1.812 & 60 & circular block bootstrap & 8 & 20260625 & Gross before costs \\
Underlying Markowitz & Sortino & 1.856 & 0.172 & 4.958 & 60 & circular block bootstrap & 8 & 20260625 & Gross before costs \\
Underlying Markowitz & Calmar & 0.590 & 0.045 & 2.544 & 60 & circular block bootstrap & 8 & 20260625 & Gross before costs \\
Underlying Markowitz & Omega & 1.929 & 1.074 & 3.877 & 60 & circular block bootstrap & 8 & 20260625 & Gross before costs \\
Underlying Markowitz & Information Ratio & 0.072 & -0.796 & 0.866 & 60 & circular block bootstrap & 8 & 20260625 & Gross before costs \\
Equity-option Greek Markowitz & Sharpe & -1.443 & -2.209 & -1.164 & 60 & circular block bootstrap & 8 & 20260625 & Post-cost research \\
Equity-option Greek Markowitz & Sortino & -1.347 & -1.895 & -1.123 & 60 & circular block bootstrap & 8 & 20260625 & Post-cost research \\
Equity-option Greek Markowitz & Calmar & -6.357 & -11.009 & -1.856 & 60 & circular block bootstrap & 8 & 20260625 & Post-cost research \\
Equity-option Greek Markowitz & Omega & 0.098 & 0.045 & 0.218 & 60 & circular block bootstrap & 8 & 20260625 & Post-cost research \\
Equity-option Greek Markowitz & Information Ratio & -1.545 & -2.532 & -1.238 & 60 & circular block bootstrap & 8 & 20260625 & Post-cost research \\
Greek Markowitz + VIX & Sharpe & -1.259 & -1.871 & -0.994 & 60 & circular block bootstrap & 8 & 20260625 & Post-cost research \\
Greek Markowitz + VIX & Sortino & -1.198 & -1.683 & -0.971 & 60 & circular block bootstrap & 8 & 20260625 & Post-cost research \\
Greek Markowitz + VIX & Calmar & -5.136 & -8.116 & -1.087 & 60 & circular block bootstrap & 8 & 20260625 & Post-cost research \\
Greek Markowitz + VIX & Omega & 0.116 & 0.054 & 0.269 & 60 & circular block bootstrap & 8 & 20260625 & Post-cost research \\
Greek Markowitz + VIX & Information Ratio & -1.345 & -2.133 & -1.069 & 60 & circular block bootstrap & 8 & 20260625 & Post-cost research \\
Beta/delta-neutral + VIX & Sharpe & -1.371 & -1.961 & -1.094 & 60 & circular block bootstrap & 8 & 20260625 & Post-cost research \\
Beta/delta-neutral + VIX & Sortino & -1.300 & -1.790 & -1.073 & 60 & circular block bootstrap & 8 & 20260625 & Post-cost research \\
Beta/delta-neutral + VIX & Calmar & -3.895 & -6.033 & -1.289 & 60 & circular block bootstrap & 8 & 20260625 & Post-cost research \\
Beta/delta-neutral + VIX & Omega & 0.151 & 0.075 & 0.321 & 60 & circular block bootstrap & 8 & 20260625 & Post-cost research \\
Beta/delta-neutral + VIX & Information Ratio & -1.508 & -2.254 & -1.234 & 60 & circular block bootstrap & 8 & 20260625 & Post-cost research \\
Cost-aware Sortino + VIX & Sharpe & -1.057 & -2.019 & -0.961 & 60 & circular block bootstrap & 8 & 20260625 & Post-cost research \\
Cost-aware Sortino + VIX & Sortino & -1.019 & -1.756 & -0.933 & 60 & circular block bootstrap & 8 & 20260625 & Post-cost research \\
Cost-aware Sortino + VIX & Calmar & -1.985 & -4.376 & -0.441 & 60 & circular block bootstrap & 8 & 20260625 & Post-cost research \\
Cost-aware Sortino + VIX & Omega & 0.016 & 0.007 & 0.032 & 60 & circular block bootstrap & 8 & 20260625 & Post-cost research \\
Cost-aware Sortino + VIX & Information Ratio & -1.003 & -2.017 & -0.902 & 60 & circular block bootstrap & 8 & 20260625 & Post-cost research \\
Equal premium & Sharpe & -1.331 & -1.975 & -0.816 & 60 & circular block bootstrap & 8 & 20260625 & Post-cost research \\
Equal premium & Sortino & -1.502 & -1.923 & -1.066 & 60 & circular block bootstrap & 8 & 20260625 & Post-cost research \\
Equal premium & Calmar & -3.051 & -4.189 & -1.543 & 60 & circular block bootstrap & 8 & 20260625 & Post-cost research \\
Equal premium & Omega & 0.359 & 0.206 & 0.547 & 60 & circular block bootstrap & 8 & 20260625 & Post-cost research \\
Equal premium & Information Ratio & -1.372 & -2.081 & -0.833 & 60 & circular block bootstrap & 8 & 20260625 & Post-cost research \\
Equal risk & Sharpe & -0.934 & -1.576 & -0.486 & 60 & circular block bootstrap & 8 & 20260625 & Post-cost research \\
Equal risk & Sortino & -1.199 & -1.714 & -0.697 & 60 & circular block bootstrap & 8 & 20260625 & Post-cost research \\
Equal risk & Calmar & -1.276 & -1.962 & -0.685 & 60 & circular block bootstrap & 8 & 20260625 & Post-cost research \\
Equal risk & Omega & 0.505 & 0.326 & 0.702 & 60 & circular block bootstrap & 8 & 20260625 & Post-cost research \\
Equal risk & Information Ratio & -1.082 & -1.764 & -0.626 & 60 & circular block bootstrap & 8 & 20260625 & Post-cost research \\
VIX hedge sleeve & Sharpe & -3.413 & -7.714 & -2.676 & 60 & circular block bootstrap & 8 & 20260625 & Post-cost research \\
VIX hedge sleeve & Sortino & -2.445 & -3.171 & -2.131 & 60 & circular block bootstrap & 8 & 20260625 & Post-cost research \\
VIX hedge sleeve & Calmar & -12.391 & -15.461 & -7.145 & 60 & circular block bootstrap & 8 & 20260625 & Post-cost research \\
VIX hedge sleeve & Omega & 0.016 & 0.004 & 0.032 & 60 & circular block bootstrap & 8 & 20260625 & Post-cost research \\
VIX hedge sleeve & Information Ratio & -3.357 & -7.275 & -2.673 & 60 & circular block bootstrap & 8 & 20260625 & Post-cost research \\
Delta-matched equities & Sharpe & 1.176 & 0.372 & 2.182 & 60 & circular block bootstrap & 8 & 20260625 & Post-cost research \\
Delta-matched equities & Sortino & 2.317 & 0.615 & 5.886 & 60 & circular block bootstrap & 8 & 20260625 & Post-cost research \\
Delta-matched equities & Calmar & 1.225 & 0.388 & 3.716 & 60 & circular block bootstrap & 8 & 20260625 & Post-cost research \\
Delta-matched equities & Omega & 2.367 & 1.301 & 5.182 & 60 & circular block bootstrap & 8 & 20260625 & Post-cost research \\
Delta-matched equities & Information Ratio & 1.041 & 0.189 & 2.072 & 60 & circular block bootstrap & 8 & 20260625 & Post-cost research \\
Underlying Markowitz & Sharpe & 0.909 & 0.103 & 1.812 & 60 & circular block bootstrap & 8 & 20260625 & Post-cost research \\
Underlying Markowitz & Sortino & 1.856 & 0.172 & 4.958 & 60 & circular block bootstrap & 8 & 20260625 & Post-cost research \\
Underlying Markowitz & Calmar & 0.590 & 0.045 & 2.544 & 60 & circular block bootstrap & 8 & 20260625 & Post-cost research \\
Underlying Markowitz & Omega & 1.929 & 1.074 & 3.877 & 60 & circular block bootstrap & 8 & 20260625 & Post-cost research \\
Underlying Markowitz & Information Ratio & 0.072 & -0.796 & 0.866 & 60 & circular block bootstrap & 8 & 20260625 & Post-cost research \\
\bottomrule
\end{tabular}
}}
  {\begin{tabular}{llll}\toprule Strategy & Metric & Estimate & Interval \\\midrule Missing & Run the empirical pipeline & -- & -- \\\bottomrule\end{tabular}}
\caption{Full bootstrap inference ledger for gross and stress-cost strategy metrics.}
\label{tab:inference_app}
\end{table}

\begin{table}[H]
\centering
\scriptsize
\IfFileExists{tables/reality_check_inference.tex}
  {\resizebox{\textwidth}{!}{\begin{tabular}{llllllllllll}
\toprule
Variant & Sharpe & Monthly Sharpe & Probabilistic Sharpe & Deflated Sharpe & DSR trials & Reality check p & Centered max stat & Centered max p95 & N & Block size & Seed \\
\midrule
Equity-option Greek Markowitz::gross & 0.8421194565895299 & 0.24309894747589325 & 0.9786072533439618 & 3.408658092459508e-11 & 9 & 0.078 & 3.160753023196929 & 3.424388436686476 & 60 & 8 & 20260625 \\
Greek Markowitz + VIX::gross & 1.3743885399954887 & 0.3967517967687661 & 0.9994784200061717 & 9.942036443637217e-08 & 9 & 0.078 & 3.160753023196929 & 3.424388436686476 & 60 & 8 & 20260625 \\
Beta/delta-neutral + VIX::gross & 1.4135317239912606 & 0.40805146067721504 & 0.9996222330231731 & 1.6814438922155963e-07 & 9 & 0.078 & 3.160753023196929 & 3.424388436686476 & 60 & 8 & 20260625 \\
Cost-aware Sortino + VIX::gross & 1.0144053186224946 & 0.29283359188704267 & 0.9978781622127668 & 4.038713604286699e-13 & 9 & 0.078 & 3.160753023196929 & 3.424388436686476 & 60 & 8 & 20260625 \\
Equal premium::gross & -0.029556763365263826 & -0.00853230264265457 & 0.4740471815096092 & 1.4656856535628173e-15 & 9 & 0.078 & 3.160753023196929 & 3.424388436686476 & 60 & 8 & 20260625 \\
Equal risk::gross & 0.08183086586465663 & 0.02362253621748983 & 0.5729898956153834 & 2.909213396306826e-15 & 9 & 0.078 & 3.160753023196929 & 3.424388436686476 & 60 & 8 & 20260625 \\
VIX hedge sleeve::gross & -6.429184989935394 & -1.8559458423045507 & 4.582392667763317e-06 & 2.8202688354747757e-12 & 9 & 0.078 & 3.160753023196929 & 3.424388436686476 & 60 & 8 & 20260625 \\
Delta-matched equities::gross & 1.1763702732937484 & 0.3395888469764097 & 0.9953982143625639 & 7.018551285374242e-08 & 9 & 0.078 & 3.160753023196929 & 3.424388436686476 & 60 & 8 & 20260625 \\
Underlying Markowitz::gross & 0.908883330100838 & 0.2623720176478412 & 0.982650547013001 & 3.947580211258329e-10 & 9 & 0.078 & 3.160753023196929 & 3.424388436686476 & 60 & 8 & 20260625 \\
Equity-option Greek Markowitz::mid & 1.0240615920008287 & 0.2956211179042176 & 0.9958848001741835 & 0.27266130728744786 & 7 & 0.078 & 3.160753023196929 & 3.424388436686476 & 60 & 8 & 20260625 \\
Greek Markowitz + VIX::mid & 0.7711086511912625 & 0.22259989366986233 & 0.9600370155790526 & 0.13421561083076827 & 7 & 0.078 & 3.160753023196929 & 3.424388436686476 & 60 & 8 & 20260625 \\
Beta/delta-neutral + VIX::mid & 1.2645573353367627 & 0.3650462589811979 & 0.9977069013011275 & 0.5054711870627524 & 7 & 0.078 & 3.160753023196929 & 3.424388436686476 & 60 & 8 & 20260625 \\
Cost-aware Sortino + VIX::mid & 0.14216345544131048 & 0.04103905463398399 & 0.6252900904885451 & 0.006071127745085553 & 7 & 0.078 & 3.160753023196929 & 3.424388436686476 & 60 & 8 & 20260625 \\
Equal premium::mid & 0.05215826661430436 & 0.015056794635116447 & 0.5468293019458382 & 0.0032539003952041755 & 7 & 0.078 & 3.160753023196929 & 3.424388436686476 & 60 & 8 & 20260625 \\
Equal risk::mid & 0.1368325564521835 & 0.039500156650786405 & 0.6215589535720385 & 0.005580984991423783 & 7 & 0.078 & 3.160753023196929 & 3.424388436686476 & 60 & 8 & 20260625 \\
VIX hedge sleeve::mid & -1.6996724333246915 & -0.4906531684570985 & 1.410586661600098e-07 & 1.9995616984377406e-19 & 7 & 0.078 & 3.160753023196929 & 3.424388436686476 & 60 & 8 & 20260625 \\
Equity-option Greek Markowitz::half\_spread & 0.9688501246318463 & 0.2796829401302995 & 0.9934794075177362 & 0.2887758280750587 & 7 & 0.078 & 3.160753023196929 & 3.424388436686476 & 60 & 8 & 20260625 \\
Greek Markowitz + VIX::half\_spread & 0.5335726940127596 & 0.15402916926025031 & 0.8852349965736346 & 0.07082348692549895 & 7 & 0.078 & 3.160753023196929 & 3.424388436686476 & 60 & 8 & 20260625 \\
Beta/delta-neutral + VIX::half\_spread & 0.9706071354884369 & 0.28019014547581034 & 0.9852643350252832 & 0.3143355313904279 & 7 & 0.078 & 3.160753023196929 & 3.424388436686476 & 60 & 8 & 20260625 \\
Cost-aware Sortino + VIX::half\_spread & 0.09439727908111553 & 0.027250147244125143 & 0.5836065916661326 & 0.007304383187074222 & 7 & 0.078 & 3.160753023196929 & 3.424388436686476 & 60 & 8 & 20260625 \\
Equal premium::half\_spread & -0.05205224625655659 & -0.01502618919407382 & 0.4548272428986697 & 0.0034731006492324528 & 7 & 0.078 & 3.160753023196929 & 3.424388436686476 & 60 & 8 & 20260625 \\
Equal risk::half\_spread & 0.01517449190340474 & 0.004380498492623262 & 0.5134507484415902 & 0.004628787852991816 & 7 & 0.078 & 3.160753023196929 & 3.424388436686476 & 60 & 8 & 20260625 \\
VIX hedge sleeve::half\_spread & -1.741614934592496 & -0.5027609256558251 & 1.7627571606747332e-08 & 9.478226752602493e-21 & 7 & 0.078 & 3.160753023196929 & 3.424388436686476 & 60 & 8 & 20260625 \\
Equity-option Greek Markowitz::full\_spread & 0.9134753475315865 & 0.26369761889772425 & 0.9899408092589378 & 0.2987736265465256 & 7 & 0.078 & 3.160753023196929 & 3.424388436686476 & 60 & 8 & 20260625 \\
Greek Markowitz + VIX::full\_spread & 0.2943184002523529 & 0.08496240380657799 & 0.7446656012700859 & 0.032333086899338205 & 7 & 0.078 & 3.160753023196929 & 3.424388436686476 & 60 & 8 & 20260625 \\
Beta/delta-neutral + VIX::full\_spread & 0.677919250357419 & 0.19569843084134259 & 0.9355215240422323 & 0.16054099931371685 & 7 & 0.078 & 3.160753023196929 & 3.424388436686476 & 60 & 8 & 20260625 \\
Cost-aware Sortino + VIX::full\_spread & 0.04656344262125476 & 0.013441708065888566 & 0.5412897255880229 & 0.008369236383558824 & 7 & 0.078 & 3.160753023196929 & 3.424388436686476 & 60 & 8 & 20260625 \\
Equal premium::full\_spread & -0.15639381474734904 & -0.04514700552198722 & 0.37082218454392535 & 0.003544055760555844 & 7 & 0.078 & 3.160753023196929 & 3.424388436686476 & 60 & 8 & 20260625 \\
Equal risk::full\_spread & -0.10651899845248258 & -0.030749386215175076 & 0.4080645087724906 & 0.003685898263592443 & 7 & 0.078 & 3.160753023196929 & 3.424388436686476 & 60 & 8 & 20260625 \\
VIX hedge sleeve::full\_spread & -1.7767478772756298 & -0.5129029326135905 & 2.5459797323930304e-09 & 7.7677544773028065e-22 & 7 & 0.078 & 3.160753023196929 & 3.424388436686476 & 60 & 8 & 20260625 \\
\bottomrule
\end{tabular}
}}
  {\begin{tabular}{llll}\toprule Variant & Sharpe & PSR & DSR \\\midrule Missing & -- & -- & -- \\\bottomrule\end{tabular}}
\caption{Reality-check inference across gross and executable-cost variants. Probabilistic
and deflated Sharpe ratios are computed in per-month units; the centered max statistic and
its bootstrap 95th percentile define the family-wise reality check (\(p=0.078\) for this
variant family).}
\label{tab:reality_check_app}
\end{table}

\begin{table}[H]
\centering
\scriptsize
\IfFileExists{tables/factor_regression.tex}
  {\resizebox{\textwidth}{!}{\begin{tabular}{lrrrrrrrrrrrrrrrrrrr}
\toprule
Strategy & Ann. alpha & Alpha HAC t & Alpha HAC se & HAC lags & $R^2$ & Residual ann. vol & N & Beta SPY & Beta AAPL & Beta AMZN & Beta GOOGL & Beta JPM & Beta META & Beta MSFT & Beta NVDA & Beta TSLA & Beta VX front & Beta dVIX & Beta dVVIX \\
\midrule
Equity-option Greek Markowitz & 0.405 & 1.141 & 0.355 & 3 & 0.618 & 0.634 & 58 & 0.005 & -1.331 & -0.875 & 0.631 & 0.590 & 0.126 & 2.585 & -0.089 & 1.136 & -1.368 & 0.046 & 0.003 \\
Greek Markowitz + VIX & 0.675 & 2.060 & 0.328 & 3 & 0.513 & 0.616 & 58 & 1.908 & -1.462 & -0.747 & 0.625 & -0.127 & -0.026 & 1.877 & 0.222 & 0.871 & -0.767 & 0.056 & -0.001 \\
Beta/delta-neutral + VIX & 0.845 & 2.753 & 0.307 & 3 & 0.429 & 0.601 & 58 & 0.715 & -1.475 & -0.612 & 0.472 & -0.579 & -0.170 & 1.336 & 0.370 & 0.815 & -0.522 & 0.038 & -0.001 \\
Equal premium & -0.416 & -0.566 & 0.735 & 3 & 0.253 & 1.647 & 58 & 2.522 & -1.714 & -0.069 & 1.943 & -1.879 & -0.344 & 3.001 & -0.374 & 0.404 & -0.749 & 0.032 & -0.001 \\
Equal risk & -0.183 & -0.334 & 0.547 & 3 & 0.253 & 1.115 & 58 & 2.041 & -1.795 & -0.113 & 1.331 & -1.484 & -0.316 & 1.968 & -0.082 & 0.321 & -0.379 & 0.020 & -0.003 \\
VIX hedge sleeve & -8.807 & -14.623 & 0.602 & 3 & 0.252 & 1.264 & 58 & 3.562 & -0.597 & -0.589 & 0.670 & -3.047 & 0.423 & -0.958 & -0.030 & -0.131 & -1.180 & 0.093 & -0.013 \\
Delta-matched equities & 0.000 & 2.229 & 0.000 & 3 & 1.000 & 0.000 & 58 & -0.000 & 0.088 & 0.539 & 0.230 & 0.590 & 0.011 & 0.442 & 0.779 & 0.628 & -0.000 & -0.000 & 0.000 \\
Underlying Markowitz & 0.000 & 2.700 & 0.000 & 3 & 1.000 & 0.000 & 58 & -0.000 & -0.085 & 0.159 & -0.096 & -0.050 & 0.029 & 0.280 & 0.191 & 0.111 & -0.000 & -0.000 & 0.000 \\
\bottomrule
\end{tabular}
}}
  {\begin{tabular}{llll}\toprule Strategy & Alpha & Controls & Status \\\midrule Missing & -- & -- & -- \\\bottomrule\end{tabular}}
\caption{Full factor regression ledger, including SPY, all eight underlying equities, VX
front returns, VIX changes, and VVIX changes. The delta-matched-equities and
underlying-Markowitz benchmark rows are exact linear combinations of the regressors
(\(R^2=1.000\), zero residual variance); their alpha \(t\)-statistics are numerical
artifacts of near-zero standard errors and should be ignored.}
\label{tab:factor_regression_full_app}
\end{table}

\subsection{Tail-path simulation diagnostics}

The next tables simulate alternative out-of-sample paths from the realized monthly strategy
returns. They are path-risk diagnostics only: simulated returns do not enter contract
selection, optimizer weights, expected-return forecasts, or any trading decision. The block
bootstrap preserves short-run dependence in the realized option return path. The
volatility-clustered simulation fits an EGARCH(1,1)-\(t\) model \citep{nelson1991} only
when the OOS sample is long enough; otherwise it reports a fixed-parameter GARCH(1,1)
residual-bootstrap fallback rather than pretending that the EGARCH parameters are precisely
estimated. Simulated wealth paths
are floored at zero (absorbing default), and the summary reports the defaulted-path share
directly. Table~\ref{tab:simulation_breach_app} gives drawdown-breach probabilities,
Table~\ref{tab:simulation_summary_app} summarizes path centers and left tails, and
Table~\ref{tab:simulation_assumptions_app} records the fitted model status per strategy.

\begin{table}[H]
\centering
\scriptsize
\IfFileExists{tables/drawdown_breach_rates.tex}
  {\resizebox{\textwidth}{!}{\begin{tabular}{llllrrrrr}
\toprule
Return basis & Strategy & Requested method & Simulation & Breach 10\% & Breach 25\% & Breach 50\% & Breach 75\% & Breach 90\% \\
\midrule
Gross before costs & Equity-option Greek Markowitz & circular\_block\_bootstrap & circular\_block\_bootstrap & 1.000 & 1.000 & 0.983 & 0.770 & 0.298 \\
Gross before costs & Equity-option Greek Markowitz & egarch\_or\_ewma & ewma\_residual\_fallback\_insufficient\_egarch\_obs & 1.000 & 1.000 & 0.987 & 0.708 & 0.296 \\
Gross before costs & Greek Markowitz + VIX & circular\_block\_bootstrap & circular\_block\_bootstrap & 1.000 & 0.996 & 0.795 & 0.135 & 0.005 \\
Gross before costs & Greek Markowitz + VIX & egarch\_or\_ewma & ewma\_residual\_fallback\_insufficient\_egarch\_obs & 1.000 & 0.999 & 0.884 & 0.379 & 0.136 \\
Gross before costs & Beta/delta-neutral + VIX & circular\_block\_bootstrap & circular\_block\_bootstrap & 1.000 & 0.996 & 0.585 & 0.063 & 0.001 \\
Gross before costs & Beta/delta-neutral + VIX & egarch\_or\_ewma & ewma\_residual\_fallback\_insufficient\_egarch\_obs & 1.000 & 0.996 & 0.671 & 0.176 & 0.037 \\
Gross before costs & Delta-matched equities & circular\_block\_bootstrap & circular\_block\_bootstrap & 1.000 & 1.000 & 0.902 & 0.548 & 0.175 \\
Gross before costs & Delta-matched equities & egarch\_or\_ewma & ewma\_residual\_fallback\_insufficient\_egarch\_obs & 1.000 & 1.000 & 0.896 & 0.358 & 0.069 \\
Gross before costs & Underlying Markowitz & circular\_block\_bootstrap & circular\_block\_bootstrap & 0.981 & 0.432 & 0.013 & 0.000 & 0.000 \\
Gross before costs & Underlying Markowitz & egarch\_or\_ewma & ewma\_residual\_fallback\_insufficient\_egarch\_obs & 0.981 & 0.272 & 0.002 & 0.000 & 0.000 \\
Gross before costs & Equal premium & circular\_block\_bootstrap & circular\_block\_bootstrap & 1.000 & 1.000 & 1.000 & 0.999 & 0.995 \\
Gross before costs & Equal premium & egarch\_or\_ewma & ewma\_residual\_fallback\_insufficient\_egarch\_obs & 0.986 & 0.986 & 0.986 & 0.986 & 0.985 \\
Gross before costs & Equal risk & circular\_block\_bootstrap & circular\_block\_bootstrap & 1.000 & 1.000 & 1.000 & 0.973 & 0.841 \\
Gross before costs & Equal risk & egarch\_or\_ewma & ewma\_residual\_fallback\_insufficient\_egarch\_obs & 1.000 & 1.000 & 1.000 & 0.993 & 0.937 \\
Gross before costs & VIX hedge sleeve & circular\_block\_bootstrap & circular\_block\_bootstrap & 0.437 & 0.437 & 0.437 & 0.437 & 0.437 \\
Gross before costs & VIX hedge sleeve & egarch\_or\_ewma & ewma\_residual\_fallback\_insufficient\_egarch\_obs & 1.000 & 1.000 & 1.000 & 1.000 & 1.000 \\
Full-spread post-cost & Equity-option Greek Markowitz & circular\_block\_bootstrap & circular\_block\_bootstrap & 1.000 & 1.000 & 0.986 & 0.789 & 0.343 \\
Full-spread post-cost & Equity-option Greek Markowitz & egarch\_or\_ewma & ewma\_residual\_fallback\_insufficient\_egarch\_obs & 1.000 & 1.000 & 0.993 & 0.756 & 0.359 \\
Full-spread post-cost & Greek Markowitz + VIX & circular\_block\_bootstrap & circular\_block\_bootstrap & 1.000 & 0.997 & 0.913 & 0.374 & 0.060 \\
Full-spread post-cost & Greek Markowitz + VIX & egarch\_or\_ewma & ewma\_residual\_fallback\_insufficient\_egarch\_obs & 1.000 & 1.000 & 0.948 & 0.503 & 0.181 \\
Full-spread post-cost & Beta/delta-neutral + VIX & circular\_block\_bootstrap & circular\_block\_bootstrap & 1.000 & 1.000 & 0.799 & 0.251 & 0.035 \\
Full-spread post-cost & Beta/delta-neutral + VIX & egarch\_or\_ewma & ewma\_residual\_fallback\_insufficient\_egarch\_obs & 1.000 & 1.000 & 0.850 & 0.295 & 0.068 \\
Full-spread post-cost & Delta-matched equities & circular\_block\_bootstrap & circular\_block\_bootstrap & 1.000 & 1.000 & 0.902 & 0.548 & 0.175 \\
Full-spread post-cost & Delta-matched equities & egarch\_or\_ewma & ewma\_residual\_fallback\_insufficient\_egarch\_obs & 1.000 & 1.000 & 0.896 & 0.358 & 0.069 \\
Full-spread post-cost & Underlying Markowitz & circular\_block\_bootstrap & circular\_block\_bootstrap & 0.981 & 0.432 & 0.013 & 0.000 & 0.000 \\
Full-spread post-cost & Underlying Markowitz & egarch\_or\_ewma & ewma\_residual\_fallback\_insufficient\_egarch\_obs & 0.981 & 0.272 & 0.002 & 0.000 & 0.000 \\
Full-spread post-cost & Equal premium & circular\_block\_bootstrap & circular\_block\_bootstrap & 1.000 & 1.000 & 1.000 & 1.000 & 0.999 \\
Full-spread post-cost & Equal premium & egarch\_or\_ewma & ewma\_residual\_fallback\_insufficient\_egarch\_obs & 0.986 & 0.986 & 0.986 & 0.986 & 0.986 \\
Full-spread post-cost & Equal risk & circular\_block\_bootstrap & circular\_block\_bootstrap & 1.000 & 1.000 & 1.000 & 0.989 & 0.938 \\
Full-spread post-cost & Equal risk & egarch\_or\_ewma & ewma\_residual\_fallback\_insufficient\_egarch\_obs & 1.000 & 1.000 & 1.000 & 1.000 & 0.984 \\
Full-spread post-cost & VIX hedge sleeve & circular\_block\_bootstrap & circular\_block\_bootstrap & 0.471 & 0.471 & 0.471 & 0.471 & 0.471 \\
Full-spread post-cost & VIX hedge sleeve & egarch\_or\_ewma & ewma\_residual\_fallback\_insufficient\_egarch\_obs & 1.000 & 1.000 & 1.000 & 1.000 & 1.000 \\
\bottomrule
\end{tabular}
}}
  {\begin{tabular}{llll}\toprule Basis & Strategy & Simulation & Breach \\\midrule Missing & -- & -- & -- \\\bottomrule\end{tabular}}
\caption{Drawdown-breach probabilities from block-bootstrap and volatility-clustered path
simulations. The thresholds are intentionally wide because option-only post-cost paths can
experience large drawdowns.}
\label{tab:simulation_breach_app}
\end{table}

\begin{table}[H]
\centering
\scriptsize
\IfFileExists{tables/simulation_summary.tex}
  {\resizebox{\textwidth}{!}{\begin{tabular}{lllrrrrrrr}
\toprule
Return basis & Strategy & Simulation & N paths & Ann. return p50 & Sortino p50 & Max DD p05 & Max DD p50 & Terminal wealth p05 & Defaulted path share \\
\midrule
Gross before costs & Equity-option Greek Markowitz & circular\_block\_bootstrap & 1000 & 0.525 & 1.826 & -0.977 & -0.842 & 0.128 & 0.000 \\
Gross before costs & Equity-option Greek Markowitz & garch11\_residual\_fallback\_insufficient\_egarch\_obs & 1000 & 0.181 & 1.238 & -0.996 & -0.873 & 0.023 & 0.002 \\
Gross before costs & Greek Markowitz + VIX & circular\_block\_bootstrap & 1000 & 1.382 & 3.451 & -0.817 & -0.604 & 3.033 & 0.000 \\
Gross before costs & Greek Markowitz + VIX & garch11\_residual\_fallback\_insufficient\_egarch\_obs & 1000 & 1.086 & 2.713 & -0.983 & -0.685 & 0.327 & 0.007 \\
Gross before costs & Beta/delta-neutral + VIX & circular\_block\_bootstrap & 1000 & 1.354 & 3.505 & -0.778 & -0.547 & 4.583 & 0.000 \\
Gross before costs & Beta/delta-neutral + VIX & garch11\_residual\_fallback\_insufficient\_egarch\_obs & 1000 & 1.130 & 3.130 & -0.910 & -0.554 & 1.634 & 0.002 \\
Gross before costs & Delta-matched equities & circular\_block\_bootstrap & 1000 & 1.006 & 2.410 & -0.964 & -0.787 & 0.717 & 0.000 \\
Gross before costs & Delta-matched equities & garch11\_residual\_fallback\_insufficient\_egarch\_obs & 1000 & 1.029 & 2.629 & -0.921 & -0.642 & 0.867 & 0.000 \\
Gross before costs & Underlying Markowitz & circular\_block\_bootstrap & 1000 & 0.176 & 1.943 & -0.422 & -0.235 & 1.035 & 0.000 \\
Gross before costs & Underlying Markowitz & garch11\_residual\_fallback\_insufficient\_egarch\_obs & 1000 & 0.143 & 1.656 & -0.398 & -0.201 & 0.974 & 0.000 \\
Gross before costs & Equal premium & circular\_block\_bootstrap & 1000 & -0.784 & -0.064 & -1.000 & -1.000 & 0.000 & 0.000 \\
Gross before costs & Equal premium & garch11\_residual\_fallback\_insufficient\_egarch\_obs & 1000 & -0.854 & -0.613 & -1.000 & -1.000 & 0.000 & 0.138 \\
Gross before costs & Equal risk & circular\_block\_bootstrap & 1000 & -0.450 & 0.153 & -1.000 & -0.983 & 0.001 & 0.000 \\
Gross before costs & Equal risk & garch11\_residual\_fallback\_insufficient\_egarch\_obs & 1000 & -0.601 & -0.349 & -1.000 & -0.996 & 0.000 & 0.002 \\
Gross before costs & VIX hedge sleeve & circular\_block\_bootstrap & 1000 & -1.000 & -3.084 & -1.000 & -1.000 & 0.000 & 1.000 \\
Gross before costs & VIX hedge sleeve & garch11\_residual\_fallback\_insufficient\_egarch\_obs & 1000 & -1.000 & -3.272 & -1.000 & -1.000 & 0.000 & 1.000 \\
Full-spread post-cost & Equity-option Greek Markowitz & circular\_block\_bootstrap & 1000 & 0.367 & 2.139 & -0.638 & -0.425 & 0.922 & 0.000 \\
Full-spread post-cost & Equity-option Greek Markowitz & garch11\_residual\_fallback\_insufficient\_egarch\_obs & 1000 & 0.140 & 0.922 & -0.792 & -0.516 & 0.424 & 0.000 \\
Full-spread post-cost & Greek Markowitz + VIX & circular\_block\_bootstrap & 1000 & 0.001 & 0.435 & -0.861 & -0.657 & 0.234 & 0.000 \\
Full-spread post-cost & Greek Markowitz + VIX & garch11\_residual\_fallback\_insufficient\_egarch\_obs & 1000 & -0.182 & -0.093 & -0.988 & -0.840 & 0.020 & 0.000 \\
Full-spread post-cost & Beta/delta-neutral + VIX & circular\_block\_bootstrap & 1000 & 0.232 & 1.124 & -0.731 & -0.514 & 0.556 & 0.000 \\
Full-spread post-cost & Beta/delta-neutral + VIX & garch11\_residual\_fallback\_insufficient\_egarch\_obs & 1000 & 0.080 & 0.632 & -0.930 & -0.645 & 0.171 & 0.000 \\
Full-spread post-cost & Delta-matched equities & circular\_block\_bootstrap & 1000 & 1.006 & 2.410 & -0.964 & -0.787 & 0.717 & 0.000 \\
Full-spread post-cost & Delta-matched equities & garch11\_residual\_fallback\_insufficient\_egarch\_obs & 1000 & 1.029 & 2.629 & -0.921 & -0.642 & 0.867 & 0.000 \\
Full-spread post-cost & Underlying Markowitz & circular\_block\_bootstrap & 1000 & 0.176 & 1.943 & -0.422 & -0.235 & 1.035 & 0.000 \\
Full-spread post-cost & Underlying Markowitz & garch11\_residual\_fallback\_insufficient\_egarch\_obs & 1000 & 0.143 & 1.656 & -0.398 & -0.201 & 0.974 & 0.000 \\
Full-spread post-cost & Equal premium & circular\_block\_bootstrap & 1000 & -0.520 & -0.290 & -0.999 & -0.985 & 0.002 & 0.000 \\
Full-spread post-cost & Equal premium & garch11\_residual\_fallback\_insufficient\_egarch\_obs & 1000 & -0.682 & -0.793 & -1.000 & -0.998 & 0.000 & 0.000 \\
Full-spread post-cost & Equal risk & circular\_block\_bootstrap & 1000 & -0.312 & -0.172 & -0.988 & -0.916 & 0.016 & 0.000 \\
Full-spread post-cost & Equal risk & garch11\_residual\_fallback\_insufficient\_egarch\_obs & 1000 & -0.298 & -0.407 & -0.985 & -0.907 & 0.019 & 0.000 \\
Full-spread post-cost & VIX hedge sleeve & circular\_block\_bootstrap & 1000 & -1.000 & -1.650 & -1.000 & -1.000 & 0.000 & 0.623 \\
Full-spread post-cost & VIX hedge sleeve & garch11\_residual\_fallback\_insufficient\_egarch\_obs & 1000 & -0.776 & -0.870 & -1.000 & -1.000 & 0.000 & 0.012 \\
\bottomrule
\end{tabular}
}}
  {\begin{tabular}{llll}\toprule Basis & Strategy & Simulation & Max DD \\\midrule Missing & -- & -- & -- \\\bottomrule\end{tabular}}
\caption{Tail-path simulation summary. Median return and Sortino describe the simulated
path center; the fifth-percentile maximum drawdown and terminal wealth columns describe the
left tail of realized-path resampling. Paths that lose all NAV are absorbed at zero
wealth and counted in the defaulted-path share column.}
\label{tab:simulation_summary_app}
\end{table}

\begin{table}[H]
\centering
\scriptsize
\IfFileExists{tables/simulation_assumptions.tex}
  {\resizebox{\textwidth}{!}{\begin{tabular}{lllllrllrrr}
\toprule
Return basis & Strategy & Method & Status & Reason & N obs & Source start & Source end & Block length & Path count & Lag1 autocorr \\
\midrule
Gross before costs & Equity-option Greek Markowitz & circular\_block\_bootstrap & ok & nan & 60 & 2021-01-29 & 2026-04-30 & 6.000 & 1000 & 0.096 \\
Gross before costs & Equity-option Greek Markowitz & ewma\_residual\_fallback\_insufficient\_egarch\_obs & fallback & insufficient\_egarch\_obs & 60 & 2021-01-29 & 2026-04-30 &  & 1000 & 0.096 \\
Gross before costs & Greek Markowitz + VIX & circular\_block\_bootstrap & ok & nan & 60 & 2021-01-29 & 2026-04-30 & 6.000 & 1000 & 0.003 \\
Gross before costs & Greek Markowitz + VIX & ewma\_residual\_fallback\_insufficient\_egarch\_obs & fallback & insufficient\_egarch\_obs & 60 & 2021-01-29 & 2026-04-30 &  & 1000 & 0.003 \\
Gross before costs & Beta/delta-neutral + VIX & circular\_block\_bootstrap & ok & nan & 60 & 2021-01-29 & 2026-04-30 & 6.000 & 1000 & -0.083 \\
Gross before costs & Beta/delta-neutral + VIX & ewma\_residual\_fallback\_insufficient\_egarch\_obs & fallback & insufficient\_egarch\_obs & 60 & 2021-01-29 & 2026-04-30 &  & 1000 & -0.083 \\
Gross before costs & Delta-matched equities & circular\_block\_bootstrap & ok & nan & 60 & 2021-01-29 & 2026-04-30 & 6.000 & 1000 & 0.014 \\
Gross before costs & Delta-matched equities & ewma\_residual\_fallback\_insufficient\_egarch\_obs & fallback & insufficient\_egarch\_obs & 60 & 2021-01-29 & 2026-04-30 &  & 1000 & 0.014 \\
Gross before costs & Underlying Markowitz & circular\_block\_bootstrap & ok & nan & 60 & 2021-01-29 & 2026-04-30 & 6.000 & 1000 & 0.041 \\
Gross before costs & Underlying Markowitz & ewma\_residual\_fallback\_insufficient\_egarch\_obs & fallback & insufficient\_egarch\_obs & 60 & 2021-01-29 & 2026-04-30 &  & 1000 & 0.041 \\
Gross before costs & Equal premium & circular\_block\_bootstrap & ok & nan & 60 & 2021-01-29 & 2026-04-30 & 6.000 & 1000 & -0.256 \\
Gross before costs & Equal premium & ewma\_residual\_fallback\_insufficient\_egarch\_obs & fallback & insufficient\_egarch\_obs & 60 & 2021-01-29 & 2026-04-30 &  & 1000 & -0.256 \\
Gross before costs & Equal risk & circular\_block\_bootstrap & ok & nan & 60 & 2021-01-29 & 2026-04-30 & 6.000 & 1000 & -0.188 \\
Gross before costs & Equal risk & ewma\_residual\_fallback\_insufficient\_egarch\_obs & fallback & insufficient\_egarch\_obs & 60 & 2021-01-29 & 2026-04-30 &  & 1000 & -0.188 \\
Gross before costs & VIX hedge sleeve & circular\_block\_bootstrap & ok & nan & 60 & 2021-01-29 & 2026-04-30 & 6.000 & 1000 & -0.125 \\
Gross before costs & VIX hedge sleeve & ewma\_residual\_fallback\_insufficient\_egarch\_obs & fallback & insufficient\_egarch\_obs & 60 & 2021-01-29 & 2026-04-30 &  & 1000 & -0.125 \\
Full-spread post-cost & Equity-option Greek Markowitz & circular\_block\_bootstrap & ok & nan & 60 & 2021-01-29 & 2026-04-30 & 6.000 & 1000 & 0.096 \\
Full-spread post-cost & Equity-option Greek Markowitz & ewma\_residual\_fallback\_insufficient\_egarch\_obs & fallback & insufficient\_egarch\_obs & 60 & 2021-01-29 & 2026-04-30 &  & 1000 & 0.096 \\
Full-spread post-cost & Greek Markowitz + VIX & circular\_block\_bootstrap & ok & nan & 60 & 2021-01-29 & 2026-04-30 & 6.000 & 1000 & -0.008 \\
Full-spread post-cost & Greek Markowitz + VIX & ewma\_residual\_fallback\_insufficient\_egarch\_obs & fallback & insufficient\_egarch\_obs & 60 & 2021-01-29 & 2026-04-30 &  & 1000 & -0.008 \\
Full-spread post-cost & Beta/delta-neutral + VIX & circular\_block\_bootstrap & ok & nan & 60 & 2021-01-29 & 2026-04-30 & 6.000 & 1000 & -0.096 \\
Full-spread post-cost & Beta/delta-neutral + VIX & ewma\_residual\_fallback\_insufficient\_egarch\_obs & fallback & insufficient\_egarch\_obs & 60 & 2021-01-29 & 2026-04-30 &  & 1000 & -0.096 \\
Full-spread post-cost & Delta-matched equities & circular\_block\_bootstrap & ok & nan & 60 & 2021-01-29 & 2026-04-30 & 6.000 & 1000 & 0.014 \\
Full-spread post-cost & Delta-matched equities & ewma\_residual\_fallback\_insufficient\_egarch\_obs & fallback & insufficient\_egarch\_obs & 60 & 2021-01-29 & 2026-04-30 &  & 1000 & 0.014 \\
Full-spread post-cost & Underlying Markowitz & circular\_block\_bootstrap & ok & nan & 60 & 2021-01-29 & 2026-04-30 & 6.000 & 1000 & 0.041 \\
Full-spread post-cost & Underlying Markowitz & ewma\_residual\_fallback\_insufficient\_egarch\_obs & fallback & insufficient\_egarch\_obs & 60 & 2021-01-29 & 2026-04-30 &  & 1000 & 0.041 \\
Full-spread post-cost & Equal premium & circular\_block\_bootstrap & ok & nan & 60 & 2021-01-29 & 2026-04-30 & 6.000 & 1000 & -0.256 \\
Full-spread post-cost & Equal premium & ewma\_residual\_fallback\_insufficient\_egarch\_obs & fallback & insufficient\_egarch\_obs & 60 & 2021-01-29 & 2026-04-30 &  & 1000 & -0.256 \\
Full-spread post-cost & Equal risk & circular\_block\_bootstrap & ok & nan & 60 & 2021-01-29 & 2026-04-30 & 6.000 & 1000 & -0.188 \\
Full-spread post-cost & Equal risk & ewma\_residual\_fallback\_insufficient\_egarch\_obs & fallback & insufficient\_egarch\_obs & 60 & 2021-01-29 & 2026-04-30 &  & 1000 & -0.188 \\
Full-spread post-cost & VIX hedge sleeve & circular\_block\_bootstrap & ok & nan & 60 & 2021-01-29 & 2026-04-30 & 6.000 & 1000 & -0.107 \\
Full-spread post-cost & VIX hedge sleeve & ewma\_residual\_fallback\_insufficient\_egarch\_obs & fallback & insufficient\_egarch\_obs & 60 & 2021-01-29 & 2026-04-30 &  & 1000 & -0.107 \\
\bottomrule
\end{tabular}
}}
  {\begin{tabular}{llll}\toprule Basis & Strategy & Method & Status \\\midrule Missing & -- & -- & -- \\\bottomrule\end{tabular}}
\caption{Simulation assumptions and model status. Fallback rows are deliberate diagnostics:
when the monthly OOS sample is too short for EGARCH estimation, the appendix uses a
fixed-parameter GARCH(1,1) residual bootstrap and records the reason.}
\label{tab:simulation_assumptions_app}
\end{table}

\subsection{Additional robustness results}

The final group holds the robustness cuts: liquidity tiers
(Table~\ref{tab:liquidity_tiers_app}), SPY-return regimes
(Table~\ref{tab:regime_performance_app}, Figure~\ref{fig:regime_sharpes_app}), VIX-level
regimes (Table~\ref{tab:vix_regime_app}, Figure~\ref{fig:vix_regime_sharpes_app}),
forecast ablations
(Table~\ref{tab:forecast_ablation_app}), leave-one-underlying-out tests
(Table~\ref{tab:loo_app}, Figure~\ref{fig:loo_app}), and the walk-forward rolling
diagnostic (Table~\ref{tab:rolling_app}). These are uncertainty screens, not proof that a
premium is stable.

\begin{table}[H]
\centering
\scriptsize
\IfFileExists{tables/liquidity_tier_performance.tex}
  {\resizebox{\textwidth}{!}{\begin{tabular}{lllllll}
\toprule
Liquidity tier & Strategy & Ann. return & Ann. vol & Sharpe & Calmar & Omega \\
\midrule
all\_eligible & Greek Markowitz & 1.2463192201296207 & 0.9068172382561567 & 1.374388539995489 & 2.033918581869829 & 2.784363176427202 \\
all\_eligible & Equal premium & -0.05552159864526898 & 1.878473564887012 & -0.02955676336526383 & -0.05553271091792252 & 0.9771765473237329 \\
all\_eligible & Equal risk & 0.10434553413680825 & 1.275136625211684 & 0.08183086586465663 & 0.10693488072635121 & 1.0641387065829775 \\
all\_eligible & Delta neutral & 1.1720642765798035 & 0.8291743699040248 & 1.4135317239912606 & 2.0573541792299705 & 2.947610573038661 \\
all\_eligible & VIX hedge sleeve & -9.28491238522107 & 1.444181867492721 & -6.429184989935395 & nan & 0.04042569095025844 \\
tight\_spread\_quartile & Greek Markowitz & 0.7303532638080936 & 1.0047849887103797 & 0.726875174305188 & 0.8741163613036393 & 1.8086255408376313 \\
tight\_spread\_quartile & Equal premium & 0.8328849438491032 & 1.954819126128125 & 0.42606752344334964 & 0.8353858747460982 & 1.3816341770771015 \\
tight\_spread\_quartile & Equal risk & 0.8890568650908819 & 1.24447750676842 & 0.7144017149811958 & 1.0024475264243362 & 1.7054047852220204 \\
tight\_spread\_quartile & Delta neutral & 0.616948156897824 & 1.839928133485132 & 0.3353110078974774 & 0.6186798622336975 & 1.2799856377055105 \\
top\_volume\_quartile & Greek Markowitz & 2.5261330738449486 & 2.0208973825933985 & 1.2500056141411722 & 3.0515281843231126 & 3.1037387879315723 \\
top\_volume\_quartile & Equal premium & -1.1699626268279582 & 2.3381506856559016 & -0.5003794811025015 & -1.1699626271251244 & 0.6936318778579352 \\
top\_volume\_quartile & Equal risk & -0.3863656502558109 & 2.183947197597906 & -0.17691162619717604 & -0.3863658542495001 & 0.8758611931290177 \\
top\_volume\_quartile & Delta neutral & 2.1254579463411787 & 1.5251047731299756 & 1.393647166928143 & 2.5319548491796637 & 3.1972700065811646 \\
top\_volume\_quartile & VIX hedge sleeve & -9.28491238522107 & 1.444181867492721 & -6.429184989935395 & nan & 0.04042569095025844 \\
\bottomrule
\end{tabular}
}}
  {\begin{tabular}{llll}\toprule Tier & Strategy & Sharpe & Calmar \\\midrule Missing & -- & -- & -- \\\bottomrule\end{tabular}}
\caption{Liquidity-tier performance. The tiers test whether the option allocation survives
when the universe is restricted to high-volume, tight-spread, high-open-interest, or jointly
liquid contracts. Tier membership is classified point-in-time from liquidity observed
through the training-window end (2020-12-31), so an allocator at the decision date could
have known the tier labels; an earlier draft classified tiers on end-of-sample liquidity.}
\label{tab:liquidity_tiers_app}
\end{table}

\begin{table}[H]
\centering
\scriptsize
\IfFileExists{tables/regime_performance.tex}
  {\resizebox{\textwidth}{!}{\begin{tabular}{llrrrrrrrrrr}
\toprule
Strategy & Regime & Months & Ann. return & Ann. vol & Sharpe & Sharpe CI lo & Sharpe CI hi & Sortino & Calmar & Omega & Info. ratio \\
\midrule
Equity-option Greek Markowitz & Down & 19 & -0.683 & 0.847 & -0.806 & -1.975 & 0.104 & -1.017 & -0.811 & 0.572 & -0.225 \\
Equity-option Greek Markowitz & Flat & 21 & 0.518 & 0.928 & 0.558 & -0.761 & 1.556 & 1.137 & 0.823 & 1.525 & -0.335 \\
Equity-option Greek Markowitz & Up & 20 & 2.787 & 1.175 & 2.373 & 0.685 & 4.553 & 8.069 & 6.615 & 6.928 & 1.808 \\
Greek Markowitz + VIX & Down & 19 & 0.806 & 0.746 & 1.081 & 0.153 & 1.938 & 2.583 & 1.918 & 2.235 & 1.764 \\
Greek Markowitz + VIX & Flat & 21 & 0.562 & 0.917 & 0.612 & -0.991 & 1.786 & 1.343 & 0.965 & 1.618 & 0.073 \\
Greek Markowitz + VIX & Up & 20 & 2.384 & 0.985 & 2.421 & 0.831 & 4.232 & 6.598 & 7.182 & 5.575 & 1.737 \\
Beta/delta-neutral + VIX & Down & 19 & 1.487 & 0.721 & 2.061 & 1.366 & 3.248 & 4.939 & 5.123 & 4.422 & 2.700 \\
Beta/delta-neutral + VIX & Flat & 21 & 0.649 & 0.815 & 0.796 & -0.666 & 1.860 & 2.154 & 1.757 & 1.990 & 0.151 \\
Beta/delta-neutral + VIX & Up & 20 & 1.423 & 0.951 & 1.496 & -0.182 & 3.218 & 3.419 & 3.357 & 3.019 & 0.777 \\
Equal premium & Down & 19 & -0.515 & 1.890 & -0.273 & -1.562 & 0.450 & -0.473 & -0.535 & 0.812 & -0.013 \\
Equal premium & Flat & 21 & -2.368 & 0.879 & -2.694 & -5.258 & -1.256 & -2.270 & -2.377 & 0.155 & -3.484 \\
Equal premium & Up & 20 & 2.809 & 2.352 & 1.195 & 0.314 & 2.010 & 3.236 & 3.712 & 2.597 & 0.917 \\
Equal risk & Down & 19 & 0.101 & 1.436 & 0.070 & -0.954 & 0.868 & 0.128 & 0.130 & 1.054 & 0.411 \\
Equal risk & Flat & 21 & -1.654 & 0.687 & -2.409 & -4.471 & -0.990 & -2.113 & -1.703 & 0.155 & -3.447 \\
Equal risk & Up & 20 & 1.954 & 1.419 & 1.377 & 0.302 & 2.376 & 3.665 & 3.495 & 2.843 & 0.915 \\
VIX hedge sleeve & Down & 19 & -7.743 & 1.368 & -5.662 & -9.800 & -4.100 & -2.982 &  & 0.021 & -5.203 \\
VIX hedge sleeve & Flat & 21 & -8.367 & 1.909 & -4.383 & -6.184 & -3.226 & -2.829 &  & 0.104 & -4.129 \\
VIX hedge sleeve & Up & 20 & -11.714 & 0.276 & -42.492 & -107.213 & -26.949 & -3.453 &  & 0.000 & -40.570 \\
Delta-matched equities & Down & 19 & -1.767 & 0.675 & -2.619 & -4.396 & -1.375 & -2.204 & -1.817 & 0.159 & -2.086 \\
Delta-matched equities & Flat & 21 & 1.272 & 0.707 & 1.800 & 0.747 & 2.788 & 5.921 & 3.839 & 5.091 & 1.172 \\
Delta-matched equities & Up & 20 & 3.630 & 0.698 & 5.203 & 3.511 & 8.613 & 34.716 & 26.891 & 45.819 & 4.428 \\
Underlying Markowitz & Down & 19 & -0.319 & 0.120 & -2.664 & -4.820 & -1.265 & -2.262 & -0.765 & 0.148 & 1.948 \\
Underlying Markowitz & Flat & 21 & 0.172 & 0.185 & 0.933 & -0.496 & 2.168 & 2.087 & 1.120 & 2.038 & -0.742 \\
Underlying Markowitz & Up & 20 & 0.658 & 0.172 & 3.829 & 2.463 & 6.121 & 15.942 & 17.455 & 15.149 & -0.128 \\
\bottomrule
\end{tabular}
}}
  {\begin{tabular}{llll}\toprule Regime & Strategy & Return & Sharpe \\\midrule Missing & -- & -- & -- \\\bottomrule\end{tabular}}
\caption{Performance by SPY-return tercile. Down, flat, and up regimes are descriptive
cuts of the held-out sample, not fitted optimizer inputs.}
\label{tab:regime_performance_app}
\end{table}

\begin{figure}[H]
\centering
\IfFileExists{figures/regime_sharpes.pdf}
  {\includegraphics[width=0.86\textwidth]{regime_sharpes.pdf}}
  {\fbox{\parbox{0.78\textwidth}{Run the empirical pipeline to generate the regime Sharpe figure.}}}
\caption{Sharpe ratios by SPY-return tercile. The figure visualizes
Table~\ref{tab:regime_performance_app}. Values below $-8$ are clipped and annotated
with the true Sharpe so that a single failed sleeve does not compress the axis.}
\label{fig:regime_sharpes_app}
\end{figure}

\begin{table}[H]
\centering
\scriptsize
\IfFileExists{tables/vix_regime_performance.tex}
  {\resizebox{\textwidth}{!}{\begin{tabular}{llrrrrrrrrrr}
\toprule
Strategy & Regime & Months & Ann. return & Ann. vol & Sharpe & Sharpe CI lo & Sharpe CI hi & Sortino & Calmar & Omega & Info. ratio \\
\midrule
Equity-option Greek Markowitz & Low VIX & 20 & 1.585 & 0.919 & 1.724 & 1.033 & 2.559 & 4.862 & 5.911 & 3.697 & 1.954 \\
Equity-option Greek Markowitz & Mid VIX & 20 & 1.104 & 1.216 & 0.907 & -0.388 & 1.845 & 2.157 & 2.351 & 2.017 & 0.841 \\
Equity-option Greek Markowitz & High VIX & 20 & -0.006 & 1.029 & -0.006 & -1.030 & 1.097 & -0.009 & -0.007 & 0.996 & -0.808 \\
Greek Markowitz + VIX & Low VIX & 20 & 1.386 & 0.846 & 1.638 & 0.265 & 3.472 & 4.774 & 4.419 & 3.257 & 2.330 \\
Greek Markowitz + VIX & Mid VIX & 20 & 1.533 & 1.054 & 1.455 & 0.412 & 2.331 & 3.812 & 4.985 & 3.198 & 1.408 \\
Greek Markowitz + VIX & High VIX & 20 & 0.821 & 0.838 & 0.979 & -0.096 & 2.111 & 2.047 & 2.057 & 2.047 & -0.223 \\
Beta/delta-neutral + VIX & Low VIX & 20 & 1.176 & 0.761 & 1.546 & 0.252 & 3.145 & 3.811 & 5.194 & 2.948 & 2.396 \\
Beta/delta-neutral + VIX & Mid VIX & 20 & 0.974 & 0.965 & 1.009 & -0.446 & 2.152 & 2.326 & 2.350 & 2.250 & 1.028 \\
Beta/delta-neutral + VIX & High VIX & 20 & 1.366 & 0.786 & 1.737 & 0.804 & 2.874 & 4.716 & 4.063 & 4.231 & 0.383 \\
Equal premium & Low VIX & 20 & 0.249 & 1.592 & 0.157 & -1.511 & 1.123 & 0.283 & 0.286 & 1.132 & 0.551 \\
Equal premium & Mid VIX & 20 & 0.944 & 2.287 & 0.413 & -0.680 & 1.136 & 0.923 & 1.004 & 1.371 & 0.398 \\
Equal premium & High VIX & 20 & -1.360 & 1.717 & -0.792 & -3.832 & 0.263 & -1.240 & -1.372 & 0.524 & -1.345 \\
Equal risk & Low VIX & 20 & 0.598 & 1.122 & 0.533 & -0.926 & 1.641 & 0.983 & 0.824 & 1.488 & 1.038 \\
Equal risk & Mid VIX & 20 & 0.348 & 1.456 & 0.239 & -0.779 & 0.996 & 0.488 & 0.436 & 1.204 & 0.231 \\
Equal risk & High VIX & 20 & -0.633 & 1.264 & -0.501 & -2.543 & 0.500 & -0.787 & -0.679 & 0.676 & -1.227 \\
VIX hedge sleeve & Low VIX & 20 & -9.290 & 1.621 & -5.730 & -20.577 & -3.438 & -3.036 & -9.290 & 0.063 & -5.189 \\
VIX hedge sleeve & Mid VIX & 20 & -10.031 & 1.331 & -7.534 & -41.311 & -4.785 & -3.203 &  & 0.038 & -7.879 \\
VIX hedge sleeve & High VIX & 20 & -8.533 & 1.406 & -6.069 & -9.076 & -4.610 & -3.031 &  & 0.018 & -7.261 \\
Delta-matched equities & Low VIX & 20 & 2.273 & 0.692 & 3.283 & 2.654 & 4.097 & 19.209 & 15.977 & 17.371 & 2.903 \\
Delta-matched equities & Mid VIX & 20 & 1.918 & 0.917 & 2.091 & 1.374 & 2.983 & 6.416 & 6.081 & 5.040 & 1.689 \\
Delta-matched equities & High VIX & 20 & -0.904 & 0.917 & -0.986 & -2.049 & -0.174 & -1.201 & -0.980 & 0.495 & -1.403 \\
Underlying Markowitz & Low VIX & 20 & 0.350 & 0.174 & 2.008 & 0.968 & 3.220 & 5.384 & 4.524 & 4.537 & 1.879 \\
Underlying Markowitz & Mid VIX & 20 & 0.386 & 0.194 & 1.994 & 1.369 & 2.695 & 7.002 & 7.409 & 5.008 & 0.746 \\
Underlying Markowitz & High VIX & 20 & -0.201 & 0.181 & -1.108 & -2.171 & -0.324 & -1.403 & -0.622 & 0.473 & -1.495 \\
\bottomrule
\end{tabular}
}}
  {\begin{tabular}{llll}\toprule Strategy & Regime & Return & Sharpe \\\midrule Missing & -- & -- & -- \\\bottomrule\end{tabular}}
\caption{Performance by VIX-level tercile.}
\label{tab:vix_regime_app}
\end{table}

\begin{figure}[H]
\centering
\IfFileExists{figures/vix_regime_sharpes.pdf}
  {\includegraphics[width=0.86\textwidth]{vix_regime_sharpes.pdf}}
  {\fbox{\parbox{0.78\textwidth}{Run the empirical pipeline to generate the VIX regime Sharpe figure.}}}
\caption{Sharpe ratios by VIX-level tercile. The figure visualizes
Table~\ref{tab:vix_regime_app}; values below $-8$ would be clipped and annotated as in
Figure~\ref{fig:regime_sharpes_app}.}
\label{fig:vix_regime_sharpes_app}
\end{figure}

The vol-of-vol regime diagnostic conditions on the prior decision date's VIX chain
feature. These features are not optimizer inputs; the table is a descriptive robustness
screen for realized returns after the feature was observable.

\begin{table}[H]
\centering
\scriptsize
\IfFileExists{tables/vol_of_vol_regime_performance.tex}
  {\resizebox{\textwidth}{!}{\begin{tabular}{lrrrr}
\toprule
Strategy & Bucket & Months & Mean monthly return & Annualized Sharpe \\
\midrule
Beta/delta-neutral + VIX & 1 & 7 & 0.101 & 1.215 \\
Beta/delta-neutral + VIX & 2 & 7 & 0.026 & 0.672 \\
Beta/delta-neutral + VIX & 3 & 7 & 0.099 & 1.238 \\
Cost-aware Sortino + VIX & 1 & 7 & 0.051 & 0.136 \\
Cost-aware Sortino + VIX & 2 & 7 & 0.642 & 1.122 \\
Cost-aware Sortino + VIX & 3 & 7 & 0.748 & 1.096 \\
Delta-matched equities & 1 & 7 & 0.097 & 1.911 \\
Delta-matched equities & 2 & 7 & 0.084 & 0.830 \\
Delta-matched equities & 3 & 7 & 0.198 & 2.224 \\
Equal premium & 1 & 7 & -0.039 & -0.273 \\
Equal premium & 2 & 7 & 0.022 & 0.133 \\
Equal premium & 3 & 7 & -0.021 & -0.147 \\
Equal risk & 1 & 7 & 0.008 & 0.069 \\
Equal risk & 2 & 7 & -0.002 & -0.023 \\
Equal risk & 3 & 7 & 0.025 & 0.236 \\
Equity-option Greek Markowitz & 1 & 7 & 0.052 & 0.558 \\
Equity-option Greek Markowitz & 2 & 7 & 0.014 & 0.205 \\
Equity-option Greek Markowitz & 3 & 7 & 0.127 & 0.890 \\
Greek Markowitz + VIX & 1 & 7 & 0.102 & 1.352 \\
Greek Markowitz + VIX & 2 & 7 & 0.030 & 0.895 \\
Greek Markowitz + VIX & 3 & 7 & 0.101 & 0.986 \\
Underlying Markowitz & 1 & 7 & 0.010 & 0.794 \\
Underlying Markowitz & 2 & 7 & 0.015 & 0.753 \\
Underlying Markowitz & 3 & 7 & 0.033 & 1.530 \\
VIX hedge sleeve & 1 & 7 & -0.876 & -12.999 \\
VIX hedge sleeve & 2 & 7 & -0.996 & -318.638 \\
VIX hedge sleeve & 3 & 7 & -0.968 & -51.079 \\
\bottomrule
\end{tabular}
}}
  {\begin{tabular}{llll}\toprule Strategy & Regime & Months & Sharpe \\\midrule Missing & -- & -- & -- \\\bottomrule\end{tabular}}
\caption{VIX chain vol-of-vol regime performance. Regime labels use the prior decision
date's diagnostic feature and do not feed expected-return forecasts.}
\label{tab:vol_of_vol_regime_app}
\end{table}

\begin{table}[H]
\centering
\scriptsize
\IfFileExists{tables/forecast_ablation_performance.tex}
  {\resizebox{0.86\textwidth}{!}{\begin{tabular}{lllllll}
\toprule
Ablation & Ann. return & Ann. vol & Sharpe & Calmar & Omega & Active assets \\
\midrule
carry\_only & 6.701470998996984 & 1.9164807271628133 & 3.496758878925927 & 2.809625893980088e-05 & 8.037348345976534 & 54 \\
variance\_risk\_premium\_only & 9.631391759962566 & 0.6700827342459892 & 14.373436693306676 & nan & nan & 54 \\
skew\_tail\_only & 4.4890452067116735 & 3.2048741838250296 & 1.4006931159319274 & 0.0030793538397452446 & 2.691677332065947 & 18 \\
vix\_regime\_only & 4.590815059335688 & 2.544245648201727 & 1.8043914362516356 & 0.0042982130732032746 & 3.3862924961758023 & 54 \\
relative\_value\_only & 2.5789118807598532 & 1.4076551542781675 & 1.8320622582327668 & 7.262838528939611 & 6.695662484712621 & 54 \\
full\_conditional\_model & 4.113956564258067 & 2.9681846972242787 & 1.3860177124776858 & 2.187721423486171 & 3.1853840590190656 & 54 \\
\bottomrule
\end{tabular}
}}
  {\begin{tabular}{llll}\toprule Ablation & Return & Sharpe & Calmar \\\midrule Missing & -- & -- & -- \\\bottomrule\end{tabular}}
\caption{Forecast ablation performance. Each row isolates one economically identified
component of the conditional option-premium forecast.}
\label{tab:forecast_ablation_app}
\end{table}

\begin{table}[H]
\centering
\scriptsize
\IfFileExists{tables/leave_one_out.tex}
  {\resizebox{\textwidth}{!}{\begin{tabular}{lrrrrrrrrrrrrr}
\toprule
Exclusion & Remaining underlyings & Option assets & Ann. return & Ann. vol & Sharpe & Sharpe CI lo & Sharpe CI hi & Sortino & Calmar & Omega & Info. ratio & Net delta & Beta SPY \\
\midrule
All underlyings & 8 & 54 & 1.246 & 0.907 & 1.374 & 0.601 & 2.209 & 3.385 & 2.034 & 2.784 & 1.200 & 3.307 & 1.297 \\
No AAPL & 7 & 48 & 1.253 & 0.949 & 1.321 & 0.571 & 2.091 & 3.233 & 2.016 & 2.673 & 1.147 & 3.216 & 1.240 \\
No AMZN & 7 & 48 & 1.286 & 0.920 & 1.397 & 0.578 & 2.242 & 3.710 & 2.097 & 2.899 & 1.223 & 3.344 & 1.178 \\
No GOOGL & 7 & 49 & 1.214 & 0.895 & 1.357 & 0.547 & 2.213 & 3.325 & 1.936 & 2.740 & 1.180 & 3.246 & 1.265 \\
No JPM & 7 & 49 & 1.202 & 0.917 & 1.311 & 0.650 & 2.015 & 3.032 & 2.540 & 2.741 & 1.106 & 3.186 & 1.105 \\
No META & 7 & 47 & 1.250 & 0.944 & 1.323 & 0.566 & 2.149 & 3.188 & 2.011 & 2.658 & 1.163 & 3.381 & 1.375 \\
No MSFT & 7 & 49 & 1.410 & 0.955 & 1.476 & 0.750 & 2.228 & 3.807 & 2.386 & 3.053 & 1.312 & 3.406 & 1.627 \\
No NVDA & 7 & 46 & 1.085 & 0.821 & 1.322 & 0.547 & 2.151 & 3.256 & 1.893 & 2.750 & 1.157 & 3.571 & 0.775 \\
No TSLA & 7 & 47 & 1.198 & 0.842 & 1.424 & 0.671 & 2.251 & 3.612 & 1.988 & 3.118 & 1.265 & 3.330 & 1.287 \\
No META/NVDA/TSLA & 5 & 32 & 0.934 & 0.837 & 1.115 & 0.336 & 1.940 & 2.279 & 1.407 & 2.352 & 0.968 & 4.019 & 1.217 \\
\bottomrule
\end{tabular}
}}
  {\begin{tabular}{llll}\toprule Exclusion & Sharpe & Delta & Status \\\midrule Missing & -- & -- & -- \\\bottomrule\end{tabular}}
\caption{Leave-one-underlying-out robustness for the VIX-enabled Greek Markowitz option
portfolio.}
\label{tab:loo_app}
\end{table}

\begin{figure}[H]
\centering
\IfFileExists{figures/leave_one_out_sharpe.pdf}
  {\includegraphics[width=0.88\textwidth]{leave_one_out_sharpe.pdf}}
  {\fbox{\parbox{0.78\textwidth}{Run the empirical pipeline to generate the leave-one-out figure.}}}
\caption{Out-of-sample Sharpe ratios for leave-one-underlying-out variants
(Table~\ref{tab:loo_app}).}
\label{fig:loo_app}
\end{figure}

\begin{table}[H]
\centering
\small
\IfFileExists{tables/rolling_oos.tex}
  {\begin{tabular}{lr}
\toprule
Diagnostic & Value \\
\midrule
Rolling 36M OOS months & 20.000 \\
Rolling 36M OOS ann. return & 1.460 \\
Rolling 36M OOS ann. vol & 0.882 \\
Rolling 36M OOS Sharpe & 1.656 \\
Rolling 36M OOS Sortino & 2.983 \\
Rolling 36M OOS Calmar & 2.561 \\
Rolling 36M OOS Omega & 3.307 \\
\bottomrule
\end{tabular}
}
  {\begin{tabular}{lll}\toprule Check & Value & Status \\\midrule Missing & -- & -- \\\bottomrule\end{tabular}}
\caption{Walk-forward rolling 36-month out-of-sample diagnostic. At each evaluation date
the representative contracts, conditional premia, and covariances are re-estimated on the
trailing 36-month window and weights come from the same constrained maximum-Sharpe solver
as the headline strategies (Section~\ref{sec:metrics}).}
\label{tab:rolling_app}
\end{table}

\subsection{Instrument and settlement convention reference}

Every candidate option is first translated into a cashflow convention. Equity options record
premium, underlying, strike, expiry, right, exercise style, mark, Greek model, payoff source,
and split adjustment. VIX options record OSI symbol, VX-forward anchor, VIX settlement
source, Black--76 Greek model, and whether expiry P\&L used exact VRO/SOQ or a settlement
proxy. Short options additionally record assignment-risk flags, simulated margin funding,
and short-premium budget use.

A weak implementation-quality score does not automatically exclude a contract. It widens
the haircut, raises the stress or capacity penalty, or downgrades the claim. This is why the
paper can show VIX option diagnostics without treating proxy settlement as exact listed
settlement.

\subsection{Proofs, solver reductions, and implementation notes}

\begin{proposition}[Greek-factor covariance identity and positive definiteness]\label{prop:cov_psd}
Under Assumption~\ref{assump:projection}, with \(\Omega_t=\Var_t(f_{t+1})\) and
\(\Sigma_{\varepsilon,t}=\Var_t(\varepsilon_{t+1})\),
\[
  \Var_t(r_{O,t+1})=B_t\Omega_tB_t^\top+\Sigma_{\varepsilon,t},
\]
and the sum is positive semidefinite. If in addition \(\Sigma_{\varepsilon,t}\succ0\),
then \(\Sigma_{O,t}\succ0\).
\end{proposition}

\begin{proof}
\(\Var_t(B_tf_{t+1}+\varepsilon_{t+1})
=B_t\Omega_tB_t^\top+\Sigma_{\varepsilon,t}
+B_t\Sigma_{f\varepsilon,t}+\Sigma_{f\varepsilon,t}^\top B_t^\top\), and
\(\Sigma_{f\varepsilon,t}=\Cov_t(f_{t+1},\varepsilon_{t+1})=0\) by the projection
construction. Both remaining terms are covariance forms, hence positive semidefinite. For
any \(x\neq0\),
\(x^\top\Sigma_{O,t}x\ge x^\top\Sigma_{\varepsilon,t}x
\ge\lambda_{\min}(\Sigma_{\varepsilon,t})\|x\|^2>0\).
\end{proof}

This is the population statement, the option-risk-model analogue of an equity factor
covariance model. As recorded in Section~\ref{sec:portfolio}, the estimated
\(\widehat\Sigma_{\varepsilon,t}\), built from pairwise-complete residual moments on an
unbalanced contract panel, need not be positive semidefinite in finite samples and is
projected onto the PSD cone before entering \(\widehat\Sigma_{O,t}\).

\begin{proposition}[Unconstrained maximum Sharpe]\label{prop:maxsharpe}
Let \(\Sigma\succ0\) and \(\mu\neq0\), and define
\(\SR(q)=q^\top\mu/\sqrt{q^\top\Sigma q}\) on the domain \(q\neq0\). Then
(i) \(\SR(cq)=\SR(q)\) for all \(c>0\);
(ii) the stationary points of \(\SR\) restricted to \(\{q:q^\top\Sigma q=1\}\) are
\(q_\pm=\pm\,\Sigma^{-1}\mu/\sqrt{\mu^\top\Sigma^{-1}\mu}\);
(iii) the maximizer is the positive root \(q_+\), unique up to positive scaling, with
optimal value \(\SR^{\ast}=\sqrt{\mu^\top\Sigma^{-1}\mu}\).
\end{proposition}

\begin{proof}
(i) is immediate. By (i) it suffices to maximize \(q^\top\mu\) on the compact ellipsoid
\(\{q:q^\top\Sigma q=1\}\), where a maximizer exists. The Lagrange condition
\(\mu=2\lambda\Sigma q\) with \(\Sigma\succ0\) gives \(q=\Sigma^{-1}\mu/(2\lambda)\), and
\(q^\top\Sigma q=1\) forces \(2\lambda=\pm\sqrt{\mu^\top\Sigma^{-1}\mu}\neq0\) because
\(\mu\neq0\). The objective at \(q_\pm\) equals \(\pm\sqrt{\mu^\top\Sigma^{-1}\mu}\), so the
maximum is attained at \(q_+\) and (iii) follows.
\end{proof}

\begin{proposition}[Conic reduction of the constrained problem]\label{prop:conic}
Let \(K\subseteq\R^n\) be a closed convex cone containing some \(q\) with
\(q^\top\mu>0\), and let \(\Sigma\succ0\). Then
\[
  \sup_{q\in K\setminus\{0\}}\SR(q)
  =\max\bigl\{\,y^\top\mu:\ y\in K,\ y^\top\Sigma y\le1\,\bigr\},
\]
and the right-hand side is a convex quadratically constrained program (the
Charnes--Cooper transformation).
\end{proposition}

\begin{proof}
For \(q\in K\setminus\{0\}\), the point \(y=q/\sqrt{q^\top\Sigma q}\) lies in \(K\)
(cone), satisfies \(y^\top\Sigma y=1\), and has \(y^\top\mu=\SR(q)\), so the right side is
at least the left. Conversely, the right side is positive by assumption, and any feasible
\(y\) with \(y^\top\mu>0\) has
\(\SR(y)=y^\top\mu/\sqrt{y^\top\Sigma y}\ge y^\top\mu\), so the left side is at least the
right. The feasible set is closed, convex, and bounded because \(\Sigma\succ0\), and the
objective is linear, so the maximum is attained.
\end{proof}

The budget functionals defining \(\mathcal K_t\) in Section~\ref{sec:portfolio} are positively homogeneous
of degree one with strictly positive limits (gross NAV, short premium, per-contract,
per-underlying, Greek, beta, and stress budgets), so they fix the admissible scale along
the optimal ray of Proposition~\ref{prop:ray}; genuinely directional restrictions, such as
a long-only sleeve, form the conic part to which Proposition~\ref{prop:conic} applies. The
implemented solver accordingly either maximizes the ratio directly over \(\mathcal K_t\)
under a gross-NAV normalization or solves the reduced convex program
\(\max\{\mu^\top y: y^\top\Sigma y\le1,\ y\in\mathcal K_t\}\) and rescales the solution to
the gross budget.

The reproducibility package exposes the model code, empirical pipeline, publication tables,
figures, strategy-return ledgers, factor regressions, attribution, regime splits, VIX
settlement records, claim-audit records, and \path{tables/empirical_summary.json}. The audit
layer independently checks data lineage, point-in-time timing, optimizer constraints, output
reproducibility, figure visibility, bibliography scope, caveat text, and LaTeX rendering.

\begin{table}[H]
\centering
\small
\IfFileExists{tables/approximation_diagnostics.tex}
  {\begin{tabular}{lr}
\toprule
Moment & Value \\
\midrule
Median residual mean & 0.076 \\
Median residual vol & 2.076 \\
Median absolute residual autocorr(1) & 0.098 \\
Systematic covariance rank & 27.000 \\
\bottomrule
\end{tabular}
}
  {\begin{tabular}{lll}\toprule Diagnostic & Value & Status \\\midrule Missing & -- & -- \\\bottomrule\end{tabular}}
\caption{Training-window residual diagnostics for the local Greek-factor approximation.}
\label{tab:approx_app}
\end{table}

The production gap remains visible. Exact VRO/SOQ settlement, executable option quotes,
calibrated costs, margin, assignment, borrow, capacity, order routing, and broker
reconciliation would be required for production claims. Those claims are not made in this
article. The current contribution is the option-only theory object, a point-in-time research
implementation, and an audit layer that prevents stronger wording when the data do not
support it.

\clearpage

\subsection{Distributional-robustness assumptions and ledgers}

Table~\ref{tab:mc_repriced_assumptions_app} records the state-model and repricing
assumptions behind Section~\ref{sec:distributional_robustness}. In particular, it
documents the one-step pricing-tenor convention, the VIX forward convention, the path
count, the block length, and the seed used by the synthetic repricing tier.

\begin{table}[H]
\centering
\scriptsize
\IfFileExists{tables/mc_repriced_assumptions.tex}
  {\resizebox{\textwidth}{!}{\begin{tabular}{llll}
\toprule
Section & Item & Value & Notes \\
\midrule
State Dimension & spot:AAPL & mu=0.019431; omega=0.00113634; alpha=0.0777105; beta=0.748476 & n\_obs=70; fallback=False; reason= \\
State Dimension & spot:AMZN & mu=0.0316915; omega=0.00203097; alpha=0; beta=0.704828 & n\_obs=70; fallback=False; reason= \\
State Dimension & spot:GOOGL & mu=0.0182634; omega=0.00348159; alpha=0.143728; beta=2.13318e-14 & n\_obs=70; fallback=False; reason= \\
State Dimension & spot:JPM & mu=0.0162303; omega=0.00362136; alpha=0.211243; beta=0 & n\_obs=70; fallback=False; reason= \\
State Dimension & spot:META & mu=0.0213947; omega=0.00167887; alpha=0.315841; beta=0.392632 & n\_obs=70; fallback=False; reason= \\
State Dimension & spot:MSFT & mu=0.0262454; omega=0.000436407; alpha=0.140304; beta=0.728673 & n\_obs=70; fallback=False; reason= \\
State Dimension & spot:NVDA & mu=0.049128; omega=0.00330517; alpha=0.0538818; beta=0.712314 & n\_obs=70; fallback=False; reason= \\
State Dimension & spot:TSLA & mu=0.0189024; omega=0.000354872; alpha=0.0604221; beta=0.939578 & n\_obs=70; fallback=False; reason= \\
State Dimension & spot:VX\_FRONT & mu=-0.00580358; omega=0.0268687; alpha=0.446917; beta=0.0984006 & n\_obs=70; fallback=False; reason= \\
State Dimension & iv:AAPL & mu=0.0081573; omega=0.0225031; alpha=0.213579; beta=0.606376 & n\_obs=69; fallback=False; reason= \\
State Dimension & iv:AMZN & mu=0.000709312; omega=0.0592597; alpha=0.042638; beta=0.453598 & n\_obs=69; fallback=False; reason= \\
State Dimension & iv:GOOGL & mu=-0.00220899; omega=0.0485703; alpha=0.00707529; beta=0.485269 & n\_obs=69; fallback=False; reason= \\
State Dimension & iv:JPM & mu=-0.00210938; omega=0.0415749; alpha=0.0733827; beta=0.415499 & n\_obs=69; fallback=False; reason= \\
State Dimension & iv:META & mu=-0.00157152; omega=0.0244162; alpha=0.205676; beta=0.547104 & n\_obs=69; fallback=False; reason= \\
State Dimension & iv:MSFT & mu=0.00613006; omega=0.0939383; alpha=0.221406; beta=1.84521e-17 & n\_obs=69; fallback=False; reason= \\
State Dimension & iv:NVDA & mu=0.000877734; omega=0.0172009; alpha=0.162236; beta=0.739677 & n\_obs=69; fallback=False; reason= \\
State Dimension & iv:TSLA & mu=-0.00981235; omega=0.0226323; alpha=0.187789; beta=0.505187 & n\_obs=69; fallback=False; reason= \\
State Dimension & iv:VX\_FRONT & mu=-0.00212937; omega=0.0459306; alpha=0.360287; beta=2.13779e-17 & n\_obs=69; fallback=False; reason= \\
State Dimension & vix:VIX & mu=0.00114242; omega=0.0338075; alpha=0.49259; beta=0.0963354 & n\_obs=68; fallback=False; reason= \\
Kernel & Simulation Method & joint\_garch\_block & joint block rows preserve empirical cross-dependence; gaussian\_copula is rank-mapped sensitivity \\
Kernel & Strike Rule & K\_t = entry state\_t * exp(train median log moneyness) & Static moneyness, repriced each month \\
Kernel & Original Tenor Reference & train median tenor\_days / 365 & Stored as tenor\_years for audit only \\
Kernel & Pricing Tenor Rule & 0.08333333333333333 & Synthetic contracts are one-step (1-month) options so premium and payoff horizons match; realized-kernel tenors average ~0.6 months \\
Kernel & Skew Freeze & contract median IV / underlying ATM median IV & 54 contracts \\
Kernel & Rate & 0.02 & Continuously-compounded annual rate used in BS and Black-76 \\
Kernel & IV Floors Caps & 0.05 to 2.0 & Applied to ATM and skew-adjusted contract IV \\
Kernel & VIX Level Caps & 5.0 to 150.0 & Applied during state reconstruction \\
Kernel & Min Mark Rule & 0.25 & Entry denominator uses max(model premium, min\_mark) in dollars \\
Kernel & Spot Anchor & contract train-end spot & State spot paths are relative to 1.0 and scaled by anchor\_spot at repricing \\
Kernel & VIX Forward Convention & simulated spot:VX\_FRONT entry forward; VIX level settlement proxy at t+1 & Premium uses the VX-front state, payoff uses simulated VIX level \\
Kernel & Seed & 20260625 &  \\
Kernel & Path Count & 1000 &  \\
Kernel & Horizon Months & 60 &  \\
Kernel & Block Length & 6 &  \\
Kernel & Min GARCH Obs & 60 &  \\
\bottomrule
\end{tabular}
}}
  {\begin{tabular}{llll}\toprule Section & Item & Value & Notes \\\midrule Missing & Run \texttt{make robustness} & -- & -- \\\bottomrule\end{tabular}}
\caption{Repriced-universe state-model and pricing assumptions.}
\label{tab:mc_repriced_assumptions_app}
\end{table}

The full fold schedule is intentionally not typeset because the blocked and CPCV ledgers
are large. The machine-readable schedule and fold ledger ship as
\path{artifacts/cv_fold_schedule.csv} and \path{artifacts/cv_fold_ledger.csv}; the split
membership, CPCV path metrics, path-month returns, PBO summary, regime tags, runtime log,
and context-consistency checks are in the adjacent \path{artifacts/cv_*.csv} files.
